% ===========================================================================
% Relatório teórico — Projeto, Parte 1
% Segurança em Redes e Sistemas de Informação (SRSI)
%
% Compilar:   make            (pdflatex + bibtex + pdflatex x2)
%       ou:   latexmk -pdf main.tex
% Overleaf:   carregar a pasta report/ inteira; compilador pdfLaTeX.
%
% Antes de entregar: preencher os campos marcados a vermelho
% (capa abaixo e dados-revisao.tex).
% ===========================================================================
\documentclass[11pt, a4paper]{article}

% ----------------------------------------------------- dados da capa -------
\newcommand{\Mestrado}{\textcolor{red}{$\langle$Nome do Mestrado$\rangle$}}
\newcommand{\UnidadeCurricular}{Segurança em Redes e Sistemas de Informação}
\newcommand{\Autores}{%
  \textcolor{red}{Nome 1}, N.º \textcolor{red}{xxxxx}\\
  \textcolor{red}{Nome 2}, N.º \textcolor{red}{yyyyy}}
\newcommand{\Data}{17-10-26}
\newcommand{\Docente}{\textcolor{red}{$\langle$Nome do Docente$\rangle$}}
\newcommand{\Titulo}{Segurança em Camadas para Comunicações IoT Baseadas em MQTT:
  Revisão Sistemática da Literatura e Proposta de Arquitetura}
\newcommand{\TituloCurto}{Segurança em camadas para MQTT}
\newcommand{\Repositorio}{https://github.com/joaoaeafonso/mqtt_secure}

% ---------------------------------------------------------------------------
% Preâmbulo: pacotes, aspeto e macros do relatório.
% ---------------------------------------------------------------------------

\usepackage[T1]{fontenc}
\usepackage[utf8]{inputenc}
\usepackage[portuguese]{babel}
\usepackage{lmodern}
\usepackage{microtype}
\usepackage[a4paper, margin=2.5cm, headheight=14pt]{geometry}
\usepackage{setspace}
\onehalfspacing

% Matemática
\usepackage{amsmath, amssymb, mathtools, bm}

% Tabelas
\usepackage{array, booktabs, tabularx, longtable, multirow, makecell}
\usepackage{colortbl}
\usepackage{threeparttable}
\usepackage{pdflscape}

% Figuras
\usepackage{graphicx}
\usepackage{float}
\usepackage[font=small, labelfont=bf, labelsep=period]{caption}
\usepackage{subcaption}
\usepackage{tikz}
\usetikzlibrary{arrows.meta, positioning, shapes.geometric, shapes.multipart,
                fit, calc, backgrounds, decorations.pathreplacing}
\usepackage{pgfplots}
\pgfplotsset{compat=1.16}

% Diversos
\usepackage{xcolor}
\usepackage{enumitem}
\usepackage{pifont}
\usepackage{csquotes}
\usepackage{url}
\usepackage{fancyhdr}
\usepackage{titlesec}
\usepackage[hidelinks, unicode, pdfusetitle]{hyperref}

% ---------------------------------------------------------------------------
% Cores (inspiradas na identidade do Iscte; ajustáveis)
% ---------------------------------------------------------------------------
\definecolor{iscte}{RGB}{0, 51, 102}
\definecolor{acento}{RGB}{0, 112, 192}
\definecolor{cinza}{RGB}{110, 110, 110}
\definecolor{fundotab}{RGB}{235, 241, 248}
\definecolor{bom}{RGB}{0, 128, 0}
\definecolor{mau}{RGB}{190, 30, 45}
\definecolor{parcial}{RGB}{200, 130, 0}

% ---------------------------------------------------------------------------
% Títulos de secção e cabeçalhos
% ---------------------------------------------------------------------------
\titleformat{\section}{\Large\bfseries\color{iscte}}{\thesection}{0.8em}{}
\titleformat{\subsection}{\large\bfseries\color{iscte}}{\thesubsection}{0.8em}{}
\titleformat{\subsubsection}{\normalsize\bfseries\color{iscte}}{\thesubsubsection}{0.8em}{}

\pagestyle{fancy}
\fancyhf{}
\fancyhead[L]{\small\color{cinza}\TituloCurto}
\fancyhead[R]{\small\color{cinza}SRSI --- Projeto, Parte 1}
\fancyfoot[C]{\small\thepage}
\renewcommand{\headrulewidth}{0.4pt}

% Nomes em português de Portugal
\addto\captionsportuguese{%
  \renewcommand{\contentsname}{Índice}%
  \renewcommand{\listfigurename}{Lista de Figuras}%
  \renewcommand{\listtablename}{Lista de Tabelas}%
  \renewcommand{\refname}{Referências}%
  \renewcommand{\abstractname}{Resumo}%
  \renewcommand{\tablename}{Tabela}%
  \renewcommand{\figurename}{Figura}%
  \renewcommand{\appendixname}{Anexo}%
}

% ---------------------------------------------------------------------------
% Macros
% ---------------------------------------------------------------------------
% Marcas usadas nas tabelas comparativas
\newcommand{\Sim}{\textcolor{bom}{\ding{51}}}          % abordado
\newcommand{\Parc}{\textcolor{parcial}{\ding{108}}}    % abordado parcialmente
\newcommand{\Nr}{\textcolor{cinza}{--}}                % não reportado
\newcommand{\Nao}{\textcolor{mau}{\ding{55}}}          % ausente por desenho

% Identificadores dos estudos incluídos (E01..E27) com citação
\newcommand{\E}[1]{\textsf{E#1}}

% Colunas de tabela
\newcolumntype{L}[1]{>{\raggedright\arraybackslash}p{#1}}
\newcolumntype{C}[1]{>{\centering\arraybackslash}p{#1}}
\newcolumntype{Y}{>{\raggedright\arraybackslash}X}

% Caixas de definição/observação
\newcommand{\caixa}[2]{%
  \begin{center}\begin{tikzpicture}
    \node[draw=acento, fill=fundotab, rounded corners=2pt, inner sep=8pt,
          text width=0.9\linewidth, align=justify] {\textbf{#1.}\ #2};
  \end{tikzpicture}\end{center}}

% ===========================================================================
% Números do processo de revisão sistemática.
%
% Gerados por report/rsl/processar.py (pesquisas/contagens.txt) e pela triagem
% em report/rsl/dados/. Valores ainda por obter aparecem a vermelho no PDF.
% ===========================================================================

\newcommand{\preencher}[1]{\textcolor{red}{\textbf{[#1]}}}

% Data em que as pesquisas foram executadas
\newcommand{\dataPesquisa}{01-11-2026}

% Registos devolvidos pela string principal, por base de dados
\newcommand{\nIEEE}{1442}
\newcommand{\nScopus}{2050}
\newcommand{\nScienceDirect}{214}               % 130 + 84 (duas consultas)
\newcommand{\nTotalIdentificados}{3706}

% String secundária (revisões): IEEE 32 + Scopus 60
\newcommand{\nSecundariaTotal}{92}
\newcommand{\nSecundariaUnicos}{73}

% Triagem (diagrama PRISMA)
\newcommand{\nDuplicados}{1216}
\newcommand{\nRemovidosTipo}{40}
\newcommand{\nTriados}{2450}                     % títulos e resumos lidos
\newcommand{\nExcluidosTA}{1729}        % excluídos por título/resumo
\newcommand{\nElegiveis}{721}          % elegíveis após a triagem
\newcommand{\nTextoIntegral}{20}                 % analisados em texto integral (seleção intencional)
\newcommand{\nExcluidosTI}{1}                    % excluídos no texto integral (CI4)
\newcommand{\nIncluidos}{19}                     % estudos primários analisados (Tabela 9)
\newcommand{\nSecundarios}{5}                    % revisões relacionadas analisadas (Secção 4)



\hypersetup{pdftitle={\Titulo}, pdfsubject={SRSI -- Projeto, Parte 1}}

\begin{document}

% ------------------------------------------------------------- capa -------
\begin{titlepage}
  \thispagestyle{empty}
  \centering
  \IfFileExists{figuras/iscte-logo.png}%
    {\includegraphics[width=0.42\textwidth]{figuras/iscte-logo.png}}%
    {\begin{tikzpicture}
       \node[draw=iscte, line width=1pt, minimum width=6cm, minimum height=2.2cm,
             text=iscte, font=\Large\bfseries] {iscte};
     \end{tikzpicture}}
  \par\vspace{2.2cm}

  {\Large Mestrado em \Mestrado\par}
  \vspace{0.6cm}
  {\LARGE\bfseries\color{iscte} \UnidadeCurricular\par}
  \vspace{2.4cm}

  \begin{flushleft}
    \large
    \textbf{Autores:}\\[0.2cm]
    \Autores\\[0.8cm]
    \textbf{Data:} \Data\\[0.8cm]
    \textbf{Título do trabalho:}\\[0.2cm]
    \Titulo\\[0.8cm]
    \textbf{Docente:} \Docente\\[0.8cm]
    \textbf{Repositório:} \url{\Repositorio}
  \end{flushleft}
  \vfill
\end{titlepage}

% ---------------------------------------------------- pré-textuais --------
\pagenumbering{roman}
\section*{Resumo}
\addcontentsline{toc}{section}{Resumo}

O MQTT é um dos protocolos de mensagens mais usados na Internet das Coisas, mas deixa a
segurança à implementação. Como todas as mensagens atravessam um \emph{broker}, o
TLS não basta: termina no próprio \emph{broker}, que vê todo o conteúdo, e não
impede a repetição de mensagens nem o abuso de uma chave comprometida. Este trabalho
apresenta uma revisão sistemática da literatura sobre mecanismos de segurança para
o MQTT, conduzida segundo um protocolo explícito (perguntas de investigação, PICOC,
\emph{strings} de pesquisa, critérios de inclusão e exclusão, e extração de
dados). As pesquisas no IEEE Xplore, no Scopus e no ScienceDirect devolveram
\nTriados{} registos únicos, triados por título e resumo; \nIncluidos{} estudos
primários (2015--2026) e cinco revisões relacionadas foram analisados em texto
integral, classificados numa taxonomia de oito categorias e comparados pelas
propriedades de segurança que garantem. A literatura trata sobretudo
confidencialidade, autenticação e integridade. A validade temporal, a deteção de
duplicados, a ligação ao tópico e o não repúdio existem, mas isolados em poucos
estudos, e a maioria das soluções confia no \emph{broker}. Daqui resulta a
proposta: propõe-se uma arquitetura em
camadas comutáveis, com defesa no transporte, um envelope fim-a-fim com AES-256-GCM
(ou ChaCha20-Poly1305) ou Ed25519 e limitação de tráfego antes de qualquer operação criptográfica, e
planeia-se a segunda parte do projeto: medir, num Raspberry Pi, o custo e os
ataques impedidos por cada camada, para escolher a configuração com a maior
proteção e o menor impacto no desempenho.

\vspace{0.6cm}
\noindent\textbf{Palavras-chave:} MQTT; Internet das Coisas; segurança fim-a-fim;
cifra autenticada; revisão sistemática da literatura; defesa em profundidade.


\clearpage
\tableofcontents
\clearpage
\listoffigures
\listoftables
\clearpage
\section*{Glossário}
\addcontentsline{toc}{section}{Glossário}

\begin{longtable}{@{}L{2.6cm}L{12.2cm}@{}}
\toprule
\textbf{Sigla} & \textbf{Significado} \\
\midrule
\endhead
AAD    & \emph{Additional Authenticated Data} --- dados autenticados mas não cifrados \\
ABE    & \emph{Attribute-Based Encryption} --- cifra baseada em atributos \\
ACL    & \emph{Access Control List} --- lista de controlo de acessos \\
AEAD   & \emph{Authenticated Encryption with Associated Data} --- cifra autenticada com dados associados \\
AES    & \emph{Advanced Encryption Standard} \\
AugPAKE & \emph{Augmented Password-Authenticated Key Exchange} \\
CA     & \emph{Certificate Authority} --- autoridade de certificação \\
CP-ABE / KP-ABE & ABE com política no criptograma / na chave \\
CRL    & \emph{Certificate Revocation List} --- lista de certificados revogados \\
DoS / DDoS & \emph{(Distributed) Denial of Service} --- negação de serviço (distribuída) \\
ECC    & \emph{Elliptic Curve Cryptography} --- criptografia de curvas elípticas \\
Ed25519 & Esquema de assinatura EdDSA sobre a curva Edwards25519 \\
GCM / GMAC & \emph{Galois/Counter Mode} / \emph{Galois Message Authentication Code} \\
HMAC   & \emph{Keyed-Hash Message Authentication Code} \\
IDS    & \emph{Intrusion Detection System} --- sistema de deteção de intrusões \\
IoT    & \emph{Internet of Things} --- Internet das Coisas \\
MAC    & \emph{Message Authentication Code} --- código de autenticação de mensagem \\
MitM   & \emph{Man-in-the-Middle} --- ataque de intermediário \\
ML     & \emph{Machine Learning} --- aprendizagem automática \\
MQTT   & \emph{Message Queuing Telemetry Transport} \\
MQTT-SN & MQTT para redes de sensores (\emph{Sensor Networks}) \\
mTLS   & TLS mútuo (autenticação do cliente por certificado) \\
OAuth  & \emph{Open Authorization} \\
PICOC  & \emph{Population, Intervention, Comparison, Outcome, Context} \\
PRISMA & \emph{Preferred Reporting Items for Systematic Reviews and Meta-Analyses} \\
QoS    & \emph{Quality of Service} --- nível de garantia de entrega no MQTT \\
QE / QF / QG & Questão estatística / focada / geral (perguntas de investigação) \\
RFC    & \emph{Request for Comments} (documento normativo do IETF) \\
RSL    & Revisão Sistemática da Literatura \\
SRSI   & Segurança em Redes e Sistemas de Informação \\
STRIDE & \emph{Spoofing, Tampering, Repudiation, Information disclosure, Denial of service, Elevation of privilege} \\
TEE    & \emph{Trusted Execution Environment} --- ambiente de execução confiável \\
TLS    & \emph{Transport Layer Security} \\
\bottomrule
\end{longtable}
\addtocounter{table}{-1} % tabela sem legenda: não conta na numeração

\clearpage

% --------------------------------------------------------- texto ----------
\pagenumbering{arabic}
\section{Introdução}
\label{sec:introducao}

\subsection{Contexto}

A Internet das Coisas (IoT) liga sensores, atuadores e serviços de
\emph{back-end} que trocam pequenas mensagens, com frequência e, muitas vezes, a
partir de dispositivos com recursos muito limitados. Entre os protocolos da camada
de aplicação usados neste domínio~\cite{dizdarevic2019survey}, o
\emph{Message Queuing Telemetry Transport} (MQTT) impôs-se pela sua simplicidade:
normalizado pela OASIS nas versões 3.1.1~\cite{oasis-mqtt311} e
5.0~\cite{oasis-mqtt5}, segue o paradigma publicação/subscrição, no qual os
clientes nunca comunicam diretamente uns com os outros. Um \emph{publisher}
envia mensagens para um \emph{tópico} (por exemplo,
\texttt{fabrica/sala1/temperatura}), e um intermediário, o \emph{broker}, entrega-as
a todos os clientes cuja subscrição corresponda a esse tópico. Implementações
abertas como o Eclipse Mosquitto~\cite{light2017mosquitto} tornaram o protocolo
omnipresente, de casas inteligentes a ambientes industriais e médicos.

\subsection{Problema}

A especificação do MQTT trata a segurança como responsabilidade da
implementação: oferece um campo de nome de utilizador e palavra-passe no pedido de
ligação e recomenda o uso de TLS, mas não define proteção da mensagem em si.
Daqui resultam três problemas que estruturam este trabalho.

\begin{enumerate}[label=\textbf{P\arabic*.}, leftmargin=*]
  \item \textbf{O \emph{broker} é um ponto de confiança total.} O TLS protege
    apenas cada salto cliente--\emph{broker}~\cite{rfc8446}. No \emph{broker}, as
    mensagens existem em claro: quem o comprometer, ou qualquer cliente a quem
    uma lista de controlo de acessos mal configurada permita subscrever o filtro
    \texttt{\#}, vê e pode reencaminhar todo o tráfego~\cite{andy2017attack,
    segarra2020mqttz}.
  \item \textbf{Uma mensagem autêntica pode ser repetida ou desviada.} Mesmo uma
    mensagem criptograficamente autêntica pode ser capturada e repetida mais tarde,
    ou publicada num tópico diferente daquele para que foi criada. A autenticação
    não impede estes ataques: são precisos, em separado, um carimbo temporal
    autenticado (validade temporal), números de sequência (anti-repetição) e a
    ligação criptográfica ao tópico.
  \item \textbf{Comprometimento de dispositivos.} Num parque com milhares de
    dispositivos, alguns acabarão comprometidos. Se uma chave válida puder
    publicar em qualquer tópico, um único dispositivo comprometido pode fazer-se
    passar por todos os outros, e o incidente de botnets como a
    Mirai~\cite{kolias2017mirai} mostra a escala que estes ataques podem atingir.
\end{enumerate}

A literatura propõe numerosas soluções parciais, desde esquemas de autenticação
e autorização a cifras baseadas em atributos e ambientes de execução confiáveis,
mas é difícil perceber que propriedades cada uma garante, que ameaças deixa por
tratar e como se combinam. É este o problema que a revisão endereça:
\emph{que mecanismos de segurança propõe o estado da arte para comunicações
baseadas em MQTT, que propriedades garantem e que lacunas persistem?}

\subsection{Objetivos}

\begin{enumerate}[label=\textbf{O\arabic*.}, leftmargin=*]
  \item Conduzir uma revisão sistemática da literatura, segundo um protocolo
    explícito e reprodutível~\cite{kitchenham2007, page2021prisma}, sobre
    mecanismos de segurança para o MQTT.
  \item Classificar as soluções encontradas numa taxonomia e compará-las segundo
    um conjunto comum de propriedades de segurança e de avaliação.
  \item Comparar a presente revisão com as revisões já publicadas sobre o tema.
  \item Identificar as lacunas do estado da arte e, a partir delas, propor uma
    arquitetura de segurança em camadas, descrita formalmente.
  \item Definir o objetivo e o plano experimental da segunda parte do projeto:
    medir, num dispositivo restrito (Raspberry Pi), o custo de desempenho e o
    benefício de segurança de cada camada, para identificar a configuração com a
    maior proteção e o menor impacto no desempenho.
\end{enumerate}

\subsection{Contributos}

\begin{itemize}[leftmargin=*]
  \item Um protocolo de revisão completo (perguntas, PICOC, \emph{strings} por base
    de dados, critérios de inclusão, exclusão e qualidade, formulário de
    extração), reutilizável e reprodutível (Secção~\ref{sec:metodologia}).
  \item Uma comparação de \nIncluidos{} estudos primários e de \nSecundarios{}
    revisões relacionadas, com uma matriz de propriedades de segurança
    (Secções~\ref{sec:relacionados} e~\ref{sec:resultados}).
  \item A formalização de mecanismos que a literatura raramente combina: ligação
    da mensagem ao tópico, janela anti-repetição, política do recetor contra a
    redução do nível de proteção e restrição chave--tópico
    (Secção~\ref{sec:proposta}).
  \item Uma arquitetura modular cuja implementação de referência é disponibilizada
    num repositório público~\cite{mqttsecure2026repo}, com documentação técnica
    completa.
\end{itemize}

\subsection{Estrutura do documento}

A Secção~\ref{sec:enquadramento} apresenta os fundamentos: o protocolo MQTT, os
seus mecanismos nativos de segurança, as primitivas criptográficas relevantes e o
modelo de ameaças. A Secção~\ref{sec:metodologia} descreve o protocolo da revisão
sistemática. A Secção~\ref{sec:relacionados} analisa as revisões já publicadas e
justifica a presente. A Secção~\ref{sec:resultados} apresenta os resultados: a
caracterização estatística dos estudos, a taxonomia, a comparação detalhada e as
respostas às perguntas de investigação. A Secção~\ref{sec:proposta} propõe a
arquitetura em camadas e posiciona-a face aos trabalhos relacionados. A
Secção~\ref{sec:plano} define o objetivo e o plano experimental da segunda parte
do projeto. A Secção~\ref{sec:discussao} discute as lacunas e as ameaças à validade, e a
Secção~\ref{sec:conclusoes} conclui e aponta o trabalho futuro. Os anexos contêm as
\emph{strings} de pesquisa adaptadas a cada base de dados e a extração de dados por
estudo.

\section{Enquadramento teórico}
\label{sec:enquadramento}

Esta secção reúne os conceitos necessários para ler a revisão e a proposta: o
funcionamento do MQTT, os mecanismos de segurança que o protocolo e os seus
\emph{brokers} oferecem nativamente, as primitivas criptográficas usadas pelas
soluções analisadas e o modelo de ameaças que serve de referência à comparação.

% ---------------------------------------------------------------------------
\subsection{O protocolo MQTT}
\label{sec:mqtt}

O MQTT é um protocolo cliente--servidor de publicação/subscrição que funciona
sobre um transporte ordenado e fiável, normalmente TCP (porta 1883) ou TLS
(porta 8883)~\cite{oasis-mqtt311}. A Figura~\ref{fig:mqtt-arquitetura} ilustra os
papéis envolvidos. Os clientes não se conhecem entre si: o desacoplamento no
espaço, no tempo e na sincronização é precisamente o que torna o protocolo
escalável, mas também o que concentra a confiança no \emph{broker}.

\begin{figure}[ht]
  \centering
  \begin{tikzpicture}[
      node distance=1.2cm and 2.2cm,
      cli/.style={draw=iscte, fill=fundotab, rounded corners=2pt, minimum width=2.6cm,
                  minimum height=0.8cm, font=\small, align=center},
      brk/.style={draw=iscte, fill=iscte!15, rounded corners=3pt, minimum width=2.8cm,
                  minimum height=2.8cm, font=\small, align=center},
      seta/.style={-{Stealth[length=2.2mm]}, thick, iscte}]
    \node[brk] (b) {\textbf{\emph{Broker}}\\[2pt]{\footnotesize filtros, ACL,}\\
                     {\footnotesize mensagens retidas}};
    \node[cli, left=of b, yshift=0.9cm] (s1) {Sensor 1\\\footnotesize\emph{publisher}};
    \node[cli, left=of b, yshift=-0.9cm] (s2) {Sensor $n$\\\footnotesize\emph{publisher}};
    \node[cli, right=of b, yshift=0.9cm] (g) {\emph{Gateway}\\\footnotesize\emph{subscriber}};
    \node[cli, right=of b, yshift=-0.9cm] (a) {Atuador\\\footnotesize\emph{subscriber}};
    \draw[seta] (s1) -- node[above, font=\scriptsize, sloped] {PUBLISH \texttt{t}} (b);
    \draw[seta] (s2) -- (b);
    \draw[seta] (b) -- node[above, font=\scriptsize, sloped] {entrega} (g);
    \draw[seta] (b) -- (a);
    \draw[decorate, decoration={brace, amplitude=5pt, mirror}, thick, acento]
      ([yshift=-0.35cm]s2.south west) -- ([yshift=-0.35cm]b.south -| b.west)
      node[midway, below=6pt, font=\scriptsize, text=acento] {TLS (salto 1)};
    \draw[decorate, decoration={brace, amplitude=5pt, mirror}, thick, acento]
      ([yshift=-0.35cm]b.south -| b.east) -- ([yshift=-0.35cm]a.south east)
      node[midway, below=6pt, font=\scriptsize, text=acento] {TLS (salto 2)};
    \draw[decorate, decoration={brace, amplitude=6pt}, thick, bom]
      ([yshift=0.35cm]s1.north west) -- ([yshift=0.35cm]g.north east)
      node[midway, above=7pt, font=\scriptsize, text=bom]
      {proteção fim-a-fim: só esta não depende da confiança no \emph{broker}};
  \end{tikzpicture}
  \caption{Arquitetura de publicação/subscrição do MQTT. O TLS protege cada salto
    em separado e termina no \emph{broker}; só a proteção fim-a-fim da mensagem
    sobrevive a um \emph{broker} não confiável.}
  \label{fig:mqtt-arquitetura}
\end{figure}

\paragraph{Pacotes.} O protocolo define catorze tipos de pacote na versão 3.1.1
(CONNECT, CONNACK, PUBLISH, PUBACK, PUBREC, PUBREL, PUBCOMP, SUBSCRIBE, SUBACK,
UNSUBSCRIBE, UNSUBACK, PINGREQ, PINGRESP e DISCONNECT), a que a versão 5.0
acrescenta AUTH~\cite{oasis-mqtt5}. Cada pacote tem um cabeçalho fixo de dois a
cinco octetos, o que explica a eficiência do protocolo em ligações limitadas.

\paragraph{Tópicos e filtros.} Os tópicos são cadeias UTF-8 hierárquicas
separadas por \texttt{/}. As subscrições usam \emph{filtros} com dois carateres
especiais: \texttt{+} corresponde a exatamente um nível e \texttt{\#} a qualquer
número de níveis, apenas no fim do filtro. Tópicos começados por \texttt{\$}
(como \texttt{\$SYS/\dots}, com estatísticas do \emph{broker}) não são abrangidos
por filtros que comecem por um destes carateres. Formalmente, sendo um tópico
$\tau = (\tau_1, \dots, \tau_m)$ e um filtro $f = (f_1, \dots, f_k)$ sequências
de níveis, define-se a correspondência por
\begin{equation}
  \operatorname{corr}(f, \tau) \iff
  \begin{cases}
    \exists j \le \min(k, m+1):\ f_j = \texttt{\#}\ \wedge\
      \forall i < j:\ f_i \in \{\tau_i, \texttt{+}\}, \\
    k = m\ \wedge\ \forall i \le k:\ f_i \in \{\tau_i, \texttt{+}\},
  \end{cases}
  \label{eq:corr}
\end{equation}
com a restrição adicional de que, se $\tau_1$ começa por \texttt{\$}, então
$f_1 \notin \{\texttt{+}, \texttt{\#}\}$. Esta relação é usada pelas listas de
controlo de acessos e, na proposta, pela restrição chave--tópico
(Secção~\ref{sec:proposta}).

\paragraph{Qualidade de serviço.} O MQTT oferece três níveis de entrega: QoS~0
(no máximo uma vez), QoS~1 (pelo menos uma vez, com possíveis duplicados marcados
pela \emph{flag} DUP) e QoS~2 (exatamente uma vez, com um protocolo de quatro
passos). Note-se que os duplicados legítimos do QoS~1 são indistinguíveis, para um
mecanismo anti-repetição, de uma repetição maliciosa --- uma tensão que a
proposta tem de resolver.

\paragraph{Mensagens retidas e sessões.} O \emph{broker} pode reter a última
mensagem de cada tópico e entregá-la a novos subscritores, e pode manter sessões
persistentes entre ligações. Ambos os mecanismos implicam que o \emph{broker}
armazena conteúdo, o que agrava o problema P1.

% ---------------------------------------------------------------------------
\subsection{Segurança nativa do MQTT e dos \emph{brokers}}
\label{sec:seguranca-nativa}

A Tabela~\ref{tab:seguranca-nativa} resume o que o protocolo e os \emph{brokers}
comuns oferecem sem extensões.

\begin{table}[ht]
  \centering
  \caption{Mecanismos de segurança nativos do MQTT e dos \emph{brokers} e respetivas limitações.}
  \label{tab:seguranca-nativa}
  \small
  \begin{tabularx}{\textwidth}{@{}L{3.1cm}L{4.5cm}Y@{}}
    \toprule
    \textbf{Mecanismo} & \textbf{O que oferece} & \textbf{Limitação} \\
    \midrule
    Utilizador e palavra-passe (CONNECT)
      & Autenticação do cliente pelo \emph{broker}
      & Em claro sem TLS; não autentica mensagens individuais nem o \emph{broker} \\
    \addlinespace
    TLS 1.2/1.3~\cite{rfc5246, rfc8446}
      & Confidencialidade e integridade de cada salto; autenticação do servidor
      & Termina no \emph{broker}; custo de estabelecimento em dispositivos
        restritos~\cite{prantl2021performance}; um cliente só controla o seu
        próprio salto~\cite{spielvogel2021tls} \\
    \addlinespace
    TLS mútuo (certificado de cliente)
      & Autenticação forte do cliente; identidade = nome comum do certificado
      & Exige infraestrutura de chave pública e revogação (CRL) \\
    \addlinespace
    Listas de controlo de acessos
      & Autorização por tópico (leitura/escrita) no \emph{broker}
      & Baseiam-se na identidade declarada se não houver autenticação; não
        protegem contra o próprio \emph{broker} \\
    \addlinespace
    Autenticação estendida (AUTH, MQTT 5)
      & Métodos de desafio--resposta definidos pela aplicação
      & Só autenticação da ligação; não protege o conteúdo \\
    \bottomrule
  \end{tabularx}
\end{table}

Nenhum destes mecanismos protege a \emph{mensagem} de ponta a ponta. É neste
espaço que se situa a maior parte das propostas da literatura e a arquitetura
proposta neste trabalho.

% ---------------------------------------------------------------------------
\subsection{Primitivas criptográficas}
\label{sec:primitivas}

\subsubsection{Cifra autenticada com dados associados}

Uma cifra AEAD é um par de algoritmos que garante, numa só operação,
confidencialidade do texto em claro e integridade e autenticidade tanto do texto
como de dados associados que seguem em claro:
\begin{align}
  (C, T) &= \operatorname{AEAD.Enc}(K, N, P, A), \label{eq:aead-enc}\\
  \operatorname{AEAD.Dec}(K, N, C, T, A) &=
    \begin{cases}
      P & \text{se } T \text{ é válido para } (N, C, A), \\
      \bot & \text{caso contrário,}
    \end{cases} \label{eq:aead-dec}
\end{align}
onde $K$ é a chave, $N$ o \emph{nonce}, $P$ o texto em claro, $A$ os dados
associados (AAD), $C$ o criptograma e $T$ a etiqueta de autenticação. O facto de
$A$ ser autenticado sem ser cifrado é o que permite ligar criptograficamente a
uma mensagem metadados que têm de permanecer legíveis, como o cabeçalho ou o
tópico MQTT.

\paragraph{AES-GCM.} No modo GCM~\cite{nist-sp800-38d}, com chave $K$ e
\emph{nonce} de 96 bits, calcula-se a subchave $H = E_K(0^{128})$ e o bloco
inicial $J_0 = N \,\|\, 0^{31} \,\|\, 1$. O texto é cifrado em modo contador a
partir de $J_0 + 1$ e a etiqueta é
\begin{equation}
  T = \operatorname{MSB}_t\!\left(
        \operatorname{GHASH}_H\!\left(A \,\|\, 0^{v} \,\|\, C \,\|\, 0^{u}
        \,\|\, [\operatorname{len}(A)]_{64} \,\|\, [\operatorname{len}(C)]_{64}\right)
        \oplus E_K(J_0)\right),
  \label{eq:gcm}
\end{equation}
onde $\operatorname{GHASH}_H$ é uma função polinomial em $\mathrm{GF}(2^{128})$ e
$u, v$ completam blocos de 128 bits. Com $P$ vazio obtém-se o GMAC, um código de
autenticação de mensagem: $T = \operatorname{AEAD.Enc}(K, N, \varnothing, A)$.

\paragraph{Limite dos \emph{nonces} aleatórios.} A segurança do GCM colapsa se o
mesmo par $(K, N)$ for usado duas vezes: o XOR dos dois criptogramas revela o XOR
dos textos e a subchave $H$ pode ser recuperada, permitindo falsificar
etiquetas~\cite{bock2016nonce}. Com \emph{nonces} aleatórios de $b = 96$ bits e
$q$ mensagens por chave, a probabilidade de colisão é limitada pelo paradoxo do
aniversário:
\begin{equation}
  \Pr[\text{colisão}] \;\le\; \frac{q(q-1)}{2^{b+1}} \;\approx\; \frac{q^{2}}{2^{97}}.
  \label{eq:aniversario}
\end{equation}
Para $q = 2^{32}$ a probabilidade é da ordem de $2^{-33}$, razão pela qual o NIST
limita a $2^{32}$ o número de invocações por chave com \emph{nonces}
aleatórios~\cite{nist-sp800-38d}. Esquemas de longa duração precisam, por isso,
de rotação de chaves.

\paragraph{ChaCha20-Poly1305.} Alternativa sem dependência de instruções AES em
\emph{hardware}, com chave de 256 bits, \emph{nonce} de 96 bits e etiqueta de 128
bits~\cite{rfc8439}; é a escolha de~\cite{sadio2019lightweight} para nós
restritos.

\subsubsection{Códigos de autenticação de mensagem}

O HMAC~\cite{rfc2104} constrói um MAC a partir de uma função de \emph{hash} $\mathcal{H}$:
\begin{equation}
  \operatorname{HMAC}(K, m) =
    \mathcal{H}\big((K' \oplus \mathit{opad}) \,\|\,
    \mathcal{H}((K' \oplus \mathit{ipad}) \,\|\, m)\big),
  \label{eq:hmac}
\end{equation}
onde $K'$ é a chave ajustada ao tamanho do bloco. Um MAC dá integridade e
autenticidade face a quem não conhece $K$, mas \emph{não} prova a origem perante
terceiros: qualquer detentor de $K$ (incluindo o recetor) pode produzir a mesma
etiqueta.

\subsubsection{Assinaturas digitais Ed25519}

O Ed25519~\cite{bernstein2012ed25519, rfc8032} é um esquema de assinatura sobre a
curva de Edwards equivalente à Curve25519, com ponto base $B$ de ordem prima
$\ell$. A partir da chave privada (32 octetos) derivam-se o escalar $s$ e a chave
pública $A = [s]B$. Para assinar $M$:
\begin{equation}
  r = \mathcal{H}(\mathit{prefixo} \,\|\, M) \bmod \ell,\quad
  R = [r]B,\quad
  k = \mathcal{H}(R \,\|\, A \,\|\, M) \bmod \ell,\quad
  S = (r + k\,s) \bmod \ell,
  \label{eq:ed25519-sign}
\end{equation}
e a assinatura é $(R, S)$, com 64 octetos. A verificação aceita sse
\begin{equation}
  [8][S]B = [8]R + [8][k]A.
  \label{eq:ed25519-verify}
\end{equation}
Como $r$ é derivado deterministicamente da chave e da mensagem, o esquema elimina
a dependência de um gerador aleatório no momento da assinatura. Ao contrário de um
MAC, só o detentor de $s$ consegue assinar, enquanto qualquer detentor de $A$
consegue verificar: a assinatura dá \emph{prova de origem} (não repúdio).

% ---------------------------------------------------------------------------
\subsection{Validade temporal e proteção contra repetição}
\label{sec:validade}

\paragraph{Janela temporal.} Uma mensagem tem \emph{validade temporal} se foi
produzida há pouco tempo, segundo um carimbo temporal autenticado pelo remetente.
Dada uma mensagem com carimbo temporal $t_m$,
recebida no instante $t$, um recetor com tolerância de idade $\Delta_a$ e de
desvio de relógio $\Delta_s$ aceita-a sse
\begin{equation}
  t - \Delta_a \;\le\; t_m \;\le\; t + \Delta_s .
  \label{eq:validade}
\end{equation}
A validade temporal responde à pergunta \enquote{esta mensagem é recente?} e
limita o tempo durante o qual uma mensagem capturada é útil, mas \emph{não}
responde a \enquote{esta mensagem já foi recebida?}: não impede a repetição
dentro da janela $\Delta_a$. É esse o papel do mecanismo seguinte.

\paragraph{Janela deslizante de números de sequência.} O mecanismo anti-repetição
do IPsec ESP~\cite{rfc4303} mantém, por remetente, o maior número de sequência
aceite $h$ e um mapa de bits $W \in \{0,1\}^{w}$ (com $w = 64$ o valor
recomendado), em que $W_i = 1$ significa que $h - i$ já foi recebido. Para uma
mensagem autêntica com número $s > 0$:
\begin{equation}
  \operatorname{aceita}(s) \iff
  \begin{cases}
    s > h, & \text{e então } W \leftarrow (W \ll (s - h)) \mid 1,\ \ h \leftarrow s; \\
    0 \le h - s < w \ \wedge\ W_{h-s} = 0, & \text{e então } W_{h-s} \leftarrow 1; \\
    \text{falso}, & \text{se } h - s \ge w \text{ ou } W_{h-s} = 1.
  \end{cases}
  \label{eq:janela}
\end{equation}
O mecanismo tolera reordenação até $w$ posições e rejeita duplicados e mensagens
antigas. Duas condições são essenciais: o estado só pode ser atualizado
\emph{depois} de verificada a autenticidade (caso contrário um atacante avança a
janela com números forjados), e cada chave deve pertencer a um único remetente.

% ---------------------------------------------------------------------------
\subsection{Limitação de tráfego}
\label{sec:token-bucket}

O balde de testemunhos (\emph{token bucket}) limita a taxa média a $\rho$
mensagens por segundo, com rajadas até $\beta$. O número de testemunhos
disponível no instante $t$, dado o estado no último acesso $t_0$, é
\begin{equation}
  b(t) = \min\!\big(\beta,\ b(t_0) + \rho\,(t - t_0)\big),
  \label{eq:token-bucket}
\end{equation}
e uma mensagem é aceite sse $b(t) \ge 1$, caso em que $b \leftarrow b(t) - 1$. Num
intervalo de duração $\Delta$, o número de mensagens aceites é limitado por
$\beta + \rho\,\Delta$, independentemente da taxa oferecida.

% ---------------------------------------------------------------------------
\subsection{Modelo de ameaças}
\label{sec:ameacas}

Considera-se um adversário que pode (i)~observar e alterar o tráfego na rede,
(ii)~ligar-se ao \emph{broker} como cliente, subscrever, publicar e repetir
mensagens, (iii)~reivindicar a identidade de outro cliente, (iv)~comprometer o
próprio \emph{broker} ou (v)~comprometer um dispositivo legítimo, obtendo as suas
chaves. A Tabela~\ref{tab:stride} organiza as ameaças segundo as categorias STRIDE,
que correspondem às classes de problemas identificadas para o MQTT
em~\cite{hintaw2023mqtt}.

\begin{table}[ht]
  \centering
  \caption{Ameaças ao MQTT segundo a classificação STRIDE e mecanismos que as mitigam.}
  \label{tab:stride}
  \small
  \begin{tabularx}{\textwidth}{@{}L{2.4cm}L{5.2cm}Y@{}}
    \toprule
    \textbf{Categoria} & \textbf{Ameaça no MQTT} & \textbf{Mitigação típica} \\
    \midrule
    \emph{Spoofing}    & Cliente reivindica a identidade de outro; mensagens injetadas em tópicos alheios
                       & Autenticação (palavras-passe, mTLS); autenticação da mensagem; restrição chave--tópico \\
    \emph{Tampering}   & Alteração do conteúdo, do cabeçalho ou do tópico em trânsito ou no \emph{broker}
                       & TLS (por salto); AEAD/MAC/assinatura com tópico nos dados associados (fim-a-fim) \\
    \emph{Repudiation} & Remetente nega ter enviado; recetor fabrica mensagens em nome do remetente
                       & Assinaturas digitais (um MAC partilhado não basta) \\
    \emph{Information disclosure} & Subscrição de \texttt{\#}; \emph{broker} curioso; tópicos \texttt{\$SYS}
                       & ACL; cifra fim-a-fim; TEE no \emph{broker}~\cite{segarra2020mqttz} \\
    \emph{Denial of service} & Inundação de mensagens; ligações lentas~\cite{vaccari2020slowite}; esgotamento de estado
                       & Limitação de tráfego; limites de pacotes e ligações; estado limitado; IDS~\cite{vaccari2020mqttset} \\
    \emph{Elevation of privilege} & Abuso de ACL baseadas em identidade não verificada; chave de um dispositivo usada para outros
                       & ACL com identidade autenticada; restrição chave--tópico; revogação \\
    \bottomrule
  \end{tabularx}
\end{table}

\paragraph{Defesa em profundidade.} Se um ataque tiver de ultrapassar $n$ camadas
de defesa independentes, cada uma contornável com probabilidade $p_i$, a
probabilidade de sucesso é
\begin{equation}
  P_{\text{sucesso}} = \prod_{i=1}^{n} p_i ,
  \label{eq:profundidade}
\end{equation}
estritamente menor do que a de qualquer camada isolada sempre que $p_i < 1$. A
hipótese de independência raramente se verifica por inteiro (por exemplo, o roubo
de uma chave de dispositivo pode afetar várias camadas), mas o modelo justifica a
combinação de mecanismos em camadas distintas: transporte, mensagem e
disponibilidade.

\section{Metodologia: protocolo da revisão sistemática}
\label{sec:metodologia}

A revisão segue as diretrizes de Kitchenham e Charters~\cite{kitchenham2007}, em
três fases (planeamento, condução e relato), e relata a seleção dos estudos com o
diagrama PRISMA 2020 (\emph{Preferred Reporting Items for Systematic Reviews and
Meta-Analyses})~\cite{page2021prisma}.

% ---------------------------------------------------------------------------
\subsection{Perguntas de investigação}
\label{sec:perguntas}

As perguntas (Tabela~\ref{tab:perguntas}) dividem-se numa pergunta geral, que
define \emph{o que} se investiga, em perguntas focadas, sobre \emph{como} os
problemas são tratados, e em perguntas estatísticas, que quantificam \emph{quanto}
é feito.

\begin{table}[ht]
  \centering
  \caption{Perguntas de investigação.}
  \label{tab:perguntas}
  \small
  \begin{tabularx}{\textwidth}{@{}L{1.1cm}Y@{}}
    \toprule
    \textbf{ID} & \textbf{Pergunta} \\
    \midrule
    QG  & Que mecanismos de segurança propõe a literatura para comunicações IoT
          baseadas em MQTT e que propriedades garantem? \\
    QF1 & Que ameaças e vulnerabilidades do MQTT são identificadas? \\
    QF2 & Que mecanismos são propostos e onde atuam (transporte, \emph{broker},
          mensagem fim-a-fim, deteção)? \\
    QF3 & Que propriedades garantem e assumem um \emph{broker} confiável? \\
    QF4 & Como são avaliadas as soluções? \\
    QF5 & Que limitações e desafios em aberto são reportados? \\
    QE1 & Quantas publicações existem por ano? \\
    QE2 & Qual a distribuição por tipo de publicação? \\
    QE3 & Quantos estudos abordam cada categoria e cada propriedade? \\
    \bottomrule
  \end{tabularx}
\end{table}

% ---------------------------------------------------------------------------
\subsection{Critérios PICOC}

\begin{table}[ht]
  \centering
  \caption{Âmbito da revisão segundo os critérios PICOC.}
  \label{tab:picoc}
  \small
  \begin{tabularx}{\textwidth}{@{}L{3cm}Y@{}}
    \toprule
    \textbf{Critério} & \textbf{Definição nesta revisão} \\
    \midrule
    \emph{Population}   & Comunicações IoT baseadas em MQTT ou MQTT-SN (MQTT para redes de sensores): clientes e
                          \emph{brokers}. \\
    \emph{Intervention} & Mecanismos de segurança: criptográficos, de autenticação e
                          autorização, de proteção do \emph{broker}, de deteção e de
                          privacidade. \\
    \emph{Comparison}   & MQTT com a segurança nativa ou apenas com TLS; mecanismos
                          propostos entre si. \\
    \emph{Outcome}      & Propriedades garantidas, ameaças mitigadas, custo e tipo de
                          avaliação. \\
    \emph{Context}      & Investigação com revisão por pares, de 2014 (normalização do
                          MQTT 3.1.1) a 2026. \\
    \bottomrule
  \end{tabularx}
\end{table}

% ---------------------------------------------------------------------------
\subsection{\emph{String} de pesquisa e bases de dados}
\label{sec:string}

A partir das perguntas identificaram-se grupos de ideias e os respetivos sinónimos
(Tabela~\ref{tab:grupos}); os termos de um grupo combinam-se com \texttt{OR} e os
grupos entre si com \texttt{AND}. A \emph{string} principal combina G1, G2 e G3 no
título, resumo e palavras-chave; a secundária acrescenta G4, para encontrar revisões
já publicadas. A \emph{string} principal inclui \texttt{"MQTT-SN"}, porque há
soluções para a variante para redes de sensores~\cite{singh2015secure,
sadio2019lightweight}, e \texttt{"intrusion"}, porque a deteção de intrusões é uma
linha própria~\cite{vaccari2020mqttset, alotaibi2024secure}.

\begin{table}[ht]
  \centering
  \caption{Grupos de ideias e sinónimos usados na \emph{string} de pesquisa.}
  \label{tab:grupos}
  \small
  \begin{tabularx}{\textwidth}{@{}L{2.8cm}Y@{}}
    \toprule
    \textbf{Grupo} & \textbf{Termos (combinados com OR)} \\
    \midrule
    G1 --- Protocolo   & \texttt{"MQTT"}, \texttt{"MQTT-SN"}, \texttt{"MQ Telemetry Transport"},
                         \texttt{"Message Queuing Telemetry Transport"},
                         \texttt{"Message Queue Telemetry Transport"} \\
    G2 --- Segurança   & \texttt{"security"}, \texttt{"secure"}, \texttt{"authentication"},
                         \texttt{"authorization"}, \texttt{"access control"},
                         \texttt{"encryption"}, \texttt{"confidentiality"},
                         \texttt{"integrity"}, \texttt{"privacy"}, \texttt{"attack*"},
                         \texttt{"vulnerabilit*"}, \texttt{"intrusion"} \\
    G3 --- Domínio     & \texttt{"Internet of Things"}, \texttt{"IoT"}, \texttt{"IIoT"},
                         \texttt{"publish/subscribe"}, \texttt{"publish-subscribe"},
                         \texttt{"machine-to-machine"}, \texttt{"M2M"} \\
    G4 --- Revisões \newline(só na secundária) & \texttt{"systematic literature"},
                         \texttt{"systematic review"}, \texttt{"literature review"},
                         \texttt{"survey"} \\
    \bottomrule
  \end{tabularx}
\end{table}

As pesquisas foram feitas no IEEE Xplore, no Scopus e no ScienceDirect, em
\dataPesquisa{}; a adaptação da \emph{string} a cada base consta do
Anexo~\ref{anx:strings}. O ScienceDirect, por não suportar o carácter \texttt{*} e
limitar os operadores booleanos, exigiu dividir a \emph{string} principal em duas
consultas, e a secundária, que excede esse limite, não pôde ser executada. O ACM
Digital Library e a Web of Science, previstos no protocolo, não foram pesquisados:
o primeiro exige subscrição para exportar resultados e a segunda não esteve
acessível; as atas ACM relevantes para o tema (por exemplo, ARES, SaT-CPS e ICPE)
estão indexadas no Scopus. Os registos foram exportados de cada base, com resumo, e
processados com um \emph{script} próprio, disponível no
repositório~\cite{mqttsecure2026repo}, que remove os tipos de documento não
admitidos pelo CI2 (normas, revistas de divulgação e capítulos de livro) e os
duplicados, por DOI e, na falta deste, por título normalizado. A
Tabela~\ref{tab:resultados-bases} regista os registos devolvidos por cada base.

\begin{table}[ht]
  \centering
  \caption{Registos devolvidos pela \emph{string} principal em cada base de dados.}
  \label{tab:resultados-bases}
  \small
  \begin{tabular}{@{}lr@{}}
    \toprule
    \textbf{Base de dados} & \textbf{Registos} \\
    \midrule
    IEEE Xplore          & \nIEEE \\
    Scopus               & \nScopus \\
    ScienceDirect (duas consultas) & \nScienceDirect \\
    \midrule
    \textbf{Total}       & \textbf{\nTotalIdentificados} \\
    \bottomrule
  \end{tabular}
\end{table}

% ---------------------------------------------------------------------------
\subsection{Critérios de inclusão e exclusão}
\label{sec:criterios}

Um estudo é incluído se satisfizer \emph{todos} os critérios de inclusão e
\emph{nenhum} dos de exclusão (Tabela~\ref{tab:criterios}). A versão em conferência do
protocolo baseado em AugPAKE (\emph{Augmented Password-Authenticated Key
Exchange})~\cite{calabretta2018token} foi considerada elegível na
triagem, mas analisou-se apenas a versão em revista~\cite{calabretta2018mqttauth},
que os autores apresentam como alargada a partir dela~\cite[§I]{calabretta2018mqttauth}
(CE2).

\begin{table}[ht]
  \centering
  \caption{Critérios de inclusão (CI) e de exclusão (CE).}
  \label{tab:criterios}
  \small
  \begin{tabularx}{\textwidth}{@{}L{1cm}Y@{}}
    \toprule
    \textbf{ID} & \textbf{Critério} \\
    \midrule
    CI1 & Propõe, implementa ou avalia um mecanismo de segurança para MQTT ou
          MQTT-SN, ou analisa ameaças específicas do MQTT. \\
    CI2 & É um estudo primário com revisão por pares (revista, conferência ou
          \emph{workshop}). \\
    CI3 & Foi publicado entre 2014 e 2026. \\
    CI4 & O texto integral está acessível, em inglês ou português. \\
    CE1 & O MQTT é apenas usado como transporte, sem contributo de segurança. \\
    CE2 & É uma versão preliminar de outro estudo incluído. \\
    CE3 & É literatura cinzenta (tese, relatório técnico, patente). \\
    CE4 & É uma revisão (analisada à parte, Secção~\ref{sec:relacionados}). \\
    CE5 & Tem menos de quatro páginas ou não descreve o mecanismo com detalhe
          suficiente. \\
    CE6 & Trata só de outro protocolo (CoAP, AMQP, DDS). \\
    \bottomrule
  \end{tabularx}
\end{table}

% ---------------------------------------------------------------------------
\subsection{Seleção dos estudos e avaliação da qualidade}
\label{sec:selecao}

\paragraph{Triagem.} Removidos os duplicados e os tipos não admitidos, os
\nTriados{} registos foram classificados por título e resumo por um único revisor,
aplicando os critérios que não exigem o texto integral (CI1--CI3, CE1, CE3, CE4 e
CE6); em caso de dúvida, o registo seguia para a fase seguinte. Por limitação de
tempo, não houve um segundo revisor independente, pelo que não se calculou a
concordância entre revisores; esta limitação é discutida na
Secção~\ref{sec:discussao}. A triagem deixou \nElegiveis{} estudos elegíveis.

\paragraph{Seleção para análise em profundidade.} O número de estudos elegíveis é
incompatível com a leitura integral de todos no prazo desta primeira parte do
projeto. Adotou-se, por isso, um desvio ao protocolo, que se declara: analisaram-se em
texto integral \nTextoIntegral{} estudos identificados numa fase exploratória,
anterior às pesquisas formais. Todos foram devolvidos pelas pesquisas e
classificados como elegíveis na triagem, mas não foram escolhidos a partir da
lista de elegíveis segundo um critério definido \emph{a priori}. Um deles foi excluído por CI4, por o texto integral não estar acessível,
pelo que o conjunto analisado tem \nIncluidos{} estudos. Esta escolha é uma ameaça à
validade da síntese, discutida na Secção~\ref{sec:discussao}, e a lista completa
das decisões de triagem está no repositório.

\paragraph{Leitura integral.} Os estudos analisados foram lidos em texto integral,
orientando a leitura pelos itens de extração da Secção~\ref{sec:extracao}; a
secção e a página que sustentam cada classificação ficaram registadas no ficheiro
de rastreabilidade do repositório.

\paragraph{\emph{Snowballing}.} A revisão de Hintaw et al.~\cite{hintaw2023mqtt},
citada na comparação de revisões de Lakshminarayana et
al.~\cite[§I]{lakshminarayana2024securing}, não foi devolvida pelas bases e foi
acrescentada às revisões relacionadas por \emph{snowballing} retrospetivo.
A Figura~\ref{fig:prisma} resume o processo.

\begin{figure}[ht]
  \centering
  \begin{tikzpicture}[
      caixa/.style={draw=iscte, fill=fundotab, rounded corners=2pt, text width=7.4cm,
                    minimum height=1.0cm, align=center, font=\small},
      lado/.style={draw=cinza, fill=white, rounded corners=2pt, text width=5.2cm,
                   minimum height=1.0cm, align=center, font=\small},
      fase/.style={draw=iscte, fill=iscte, text=white, font=\tiny\bfseries,
                   rotate=90, minimum width=1.5cm, minimum height=0.5cm},
      seta/.style={-{Stealth[length=2mm]}, thick, iscte}]
    \node[caixa] (id) {Registos identificados nas bases de dados\\
      ($n = \nTotalIdentificados$)\\[-1pt]
      {\scriptsize IEEE \nIEEE; Scopus \nScopus; ScienceDirect \nScienceDirect}};
    \node[lado, right=0.6cm of id] (dup) {Removidos antes da triagem:\\
      duplicados ($n = \nDuplicados$)\\tipo não admitido ($n = \nRemovidosTipo$)};
    \node[caixa, below=0.9cm of id] (tri) {Registos triados por título e resumo\\($n = \nTriados$)};
    \node[lado, right=0.6cm of tri] (exta) {Excluídos (CI/CE)\\($n = \nExcluidosTA$)};
    \node[caixa, below=0.9cm of tri] (ti) {Estudos elegíveis ($n = \nElegiveis$)\\
      analisados em texto integral, da fase exploratória ($n = \nTextoIntegral$)};
    \node[lado, right=0.6cm of ti] (exti) {Excluídos no texto integral: CI4\\($n = \nExcluidosTI$)};
    \node[caixa, below=0.9cm of ti, fill=iscte!15] (inc) {\textbf{Estudos primários analisados}\\($n = \nIncluidos$)};
    \node[lado, right=0.6cm of inc] (sec) {Revisões relacionadas: \nSecundarios{}\\
      (\nSecundariaUnicos{} devolvidas pela \emph{string} secundária; 1 por \emph{snowballing})};
    \draw[seta] (id) -- (tri);
    \draw[seta] (tri) -- (ti);
    \draw[seta] (ti) -- (inc);
    \draw[seta] (id) -- (dup);
    \draw[seta] (tri) -- (exta);
    \draw[seta] (ti) -- (exti);
    \node[fase, left=0.25cm of id] {Identificação};
    \node[fase, left=0.25cm of tri] {Triagem};
    \node[fase, left=0.25cm of ti] {Elegibilidade};
    \node[fase, left=0.25cm of inc] {Inclusão};
  \end{tikzpicture}
  \caption{Diagrama de fluxo PRISMA 2020 da seleção dos estudos.}
  \label{fig:prisma}
\end{figure}

% ---------------------------------------------------------------------------
\subsection{Extração de dados}
\label{sec:extracao}

Para cada estudo analisado registaram-se os seguintes itens: ano e tipo de publicação
(QE1, QE2); tipo de contributo e categoria da taxonomia (QF1, QF2, QE3); ponto de
atuação e primitivas usadas (QF2); propriedades garantidas e confiança no
\emph{broker} (QF3, QE3); tipo de avaliação e uso de dispositivos restritos (QF4);
e limitações declaradas (QF5). O resumo da extração consta do
Anexo~\ref{anx:extracao}.

\section{Trabalhos relacionados: revisões existentes}
\label{sec:relacionados}

A \emph{string} secundária devolveu \nSecundariaTotal{} registos
(\nSecundariaUnicos{} únicos). Destes, e das revisões excluídas por CE4 na triagem
da \emph{string} principal, selecionaram-se as quatro mais diretamente
relacionadas com a segurança do MQTT; a quinta, de Hintaw et al., foi acrescentada
por \emph{snowballing} (Secção~\ref{sec:selecao}).

\begin{itemize}[leftmargin=*]
  \item \textbf{Perrone et al.\ (2017)}~\cite{perrone2017mirai}, escrito após os
    ataques da botnet Mirai, descreve as soluções de segurança habitualmente usadas
    em instalações MQTT reais e fecha com um levantamento de temas de investigação
    que os próprios autores descrevem como conciso e não exaustivo.
  \item \textbf{Nebbione e Calzarossa (2020)}~\cite{nebbione2020security} revêem a
    segurança dos protocolos da camada de aplicação da IoT, incluindo o MQTT, as
    vulnerabilidades publicadas (CVE) e boas práticas de mitigação.
  \item \textbf{Hintaw et al.\ (2023)}~\cite{hintaw2023mqtt} organizam os problemas
    de segurança do MQTT em negação de serviço, falsificação de identidade,
    divulgação de informação, elevação de privilégios e adulteração, e apontam o
    custo da segurança em dispositivos a bateria como desafio em aberto.
  \item \textbf{Lakshminarayana et al.\ (2024)}~\cite{lakshminarayana2024securing}
    classificam as vulnerabilidades e os ataques ao MQTT e discutem as defesas, com
    ênfase na identificação de vulnerabilidades e na deteção de ataques.
  \item \textbf{Laghari et al.\ (2024)}~\cite{laghari2024securing} apresentam uma
    taxonomia de ataques ao ecossistema MQTT e apontam às revisões anteriores uma
    discussão ampla mas superficial.
\end{itemize}

\begin{table}[ht]
  \centering
  \caption{Comparação das revisões relacionadas com a presente revisão.
    \Sim{} abordado; \Parc{} parcialmente; \Nr{} não reportado. $^{*}$Classificada
    a partir do resumo: o texto integral não estava acessível.}
  \label{tab:revisoes}
  \small
  \setlength{\tabcolsep}{3pt}
  \begin{tabular}{@{}lcC{1.45cm}C{1.75cm}*{4}{C{1.45cm}}@{}}
    \toprule
    \textbf{Revisão} & \textbf{Ano}
      & {\scriptsize\textbf{Só MQTT}}
      & {\scriptsize\textbf{Taxonomia de ataques}}
      & {\scriptsize\textbf{Revê defesas}}
      & {\scriptsize\textbf{Compara por propriedade}}
      & {\scriptsize\textbf{Protocolo RSL explícito}}
      & {\scriptsize\textbf{Propõe arquitetura}} \\
    \midrule
    Perrone et al.~\cite{perrone2017mirai}              & 2017 & \Sim  & \Nr   & \Sim & \Nr & \Nr & \Nr \\
    Nebbione e Calzarossa~\cite{nebbione2020security}   & 2020 & \Parc & \Sim  & \Sim & \Nr & \Nr & \Nr \\
    Hintaw et al.~\cite{hintaw2023mqtt}$^{*}$          & 2023 & \Sim  & \Parc & \Sim & \Nr & \Nr & \Nr \\
    Lakshminarayana et al.~\cite{lakshminarayana2024securing} & 2024 & \Sim & \Sim & \Sim & \Nr & \Nr & \Nr \\
    Laghari et al.~\cite{laghari2024securing}           & 2024 & \Sim  & \Sim  & \Sim & \Nr & \Parc & \Nr \\
    \midrule
    \rowcolor{fundotab}
    \textbf{Este trabalho}                              & 2026 & \Sim  & \Parc & \Sim & \Sim & \Sim & \Sim \\
    \bottomrule
  \end{tabular}
\end{table}

As revisões existentes organizam-se em torno dos \emph{ataques}
(Tabela~\ref{tab:revisoes}). Só Laghari et al.\ descrevem a pesquisa e a seleção,
e apenas das revisões que comparam, com uma única palavra-chave em três
fontes~\cite[§I]{laghari2024securing}. Esta revisão distingue-se por relatar um
protocolo explícito e reprodutível, por comparar os estudos segundo as propriedades que garantem à
\emph{mensagem} (incluindo validade temporal, anti-repetição e não repúdio) e a
confiança que depositam no \emph{broker}, por cobrir trabalhos até 2026 e por
ligar as lacunas encontradas a uma proposta concreta.

\section{Resultados da revisão}
\label{sec:resultados}

A Tabela~\ref{tab:estudos} identifica os \nIncluidos{} estudos primários
analisados em texto integral (Secção~\ref{sec:selecao}), numerados por ordem
cronológica (E01--E19), com o tipo de publicação e a categoria que lhes foi
atribuída na taxonomia da Secção~\ref{sec:taxonomia}.

\begin{table}[ht]
  \centering
  \caption{Estudos primários incluídos. Tipo: R revista, C conferência, W
    \emph{workshop}. Categorias: ver Figura~\ref{fig:taxonomia}.}
  \label{tab:estudos}
  \footnotesize
  \setlength{\tabcolsep}{4pt}
  \begin{tabular}{@{}llccc@{\hspace{1.2em}}llccc@{}}
    \toprule
    \textbf{ID} & \textbf{Estudo} & \textbf{Ano} & \textbf{Tipo} & \textbf{Cat.} &
    \textbf{ID} & \textbf{Estudo} & \textbf{Ano} & \textbf{Tipo} & \textbf{Cat.} \\
    \midrule
    \E{01} & Singh et al.~\cite{singh2015secure}               & 2015 & C & C4 &
    \E{11} & Vaccari et al.~\cite{vaccari2020mqttset}          & 2020 & R & C7 \\
    \E{02} & Niruntasukrat et al.~\cite{niruntasukrat2016authorization} & 2016 & W & C3 &
    \E{12} & Segarra et al.~\cite{segarra2020mqttz}            & 2020 & C & C6 \\
    \E{03} & Rizzardi et al.~\cite{rizzardi2016aups}           & 2016 & R & C3 &
    \E{13} & Dahlmanns et al.~\cite{dahlmanns2021entrust}      & 2021 & W & C5 \\
    \E{04} & Andy et al.~\cite{andy2017attack}                 & 2017 & C & C1 &
    \E{14} & Spielvogel et al.~\cite{spielvogel2021tls}        & 2021 & W & C6 \\
    \E{05} & Calabretta et al.~\cite{calabretta2018mqttauth}   & 2018 & R & C3 &
    \E{15} & Prantl et al.~\cite{prantl2021performance}        & 2021 & C & C6 \\
    \E{06} & Hernández Ramos et al.~\cite{hernandez2018fuzzing} & 2018 & R & C1 &
    \E{16} & Van Glabbeek et al.~\cite{vanglabbeek2022flexible} & 2022 & R & C5 \\
    \E{07} & Dinculeană e Cheng~\cite{dinculeana2019vulnerabilities} & 2019 & R & C5 &
    \E{17} & Buccafurri et al.~\cite{buccafurri2023mqtta}      & 2023 & R & C8 \\
    \E{08} & Malina et al.~\cite{malina2019secure}             & 2019 & C & C5 &
    \E{18} & Alotaibi et al.~\cite{alotaibi2024secure}         & 2024 & R & C7 \\
    \E{09} & Sadio et al.~\cite{sadio2019lightweight}          & 2019 & C & C5 &
    \E{19} & Javadi et al.~\cite{javadi2026cmap}               & 2026 & R & C2 \\
    \E{10} & Vaccari et al.~\cite{vaccari2020slowite}          & 2020 & R & C1 & & & & & \\
    \bottomrule
  \end{tabular}
\end{table}

% ---------------------------------------------------------------------------
\subsection{Caracterização dos estudos (QE1, QE2)}
\label{sec:estatistica}

\paragraph{Publicações por ano (QE1).} A Figura~\ref{fig:por-ano} mostra a
distribuição temporal dos \nElegiveis{} estudos elegíveis após a triagem. O
interesse surge logo após a normalização do MQTT~3.1.1 (2014) e cresce de forma
contínua: os anos de 2023 a 2026 reúnem 410 estudos ($56{,}9\,\%$), e o ano de 2026
está contado só até à data das pesquisas. Destes estudos elegíveis, só
\nIncluidos{} foram analisados em texto integral (Secção~\ref{sec:selecao}); as
restantes estatísticas desta secção referem-se a esse conjunto.

\begin{figure}[ht]
  \centering
  \begin{tikzpicture}
    \begin{axis}[
        ybar, bar width=12pt, width=0.92\textwidth, height=5.2cm,
        ymin=0, ymax=130, ytick={0,20,...,120},
        symbolic x coords={2014,2015,2016,2017,2018,2019,2020,2021,2022,2023,2024,2025,2026},
        xtick=data, x tick label style={font=\small},
        ylabel={N.º de estudos}, ylabel style={font=\small},
        nodes near coords, nodes near coords style={font=\scriptsize},
        axis lines*=left, ymajorgrids, grid style={cinza!25},
        enlarge x limits=0.05]
      \addplot[fill=acento!70, draw=iscte] coordinates {
        (2014,2) (2015,4) (2016,9) (2017,17) (2018,30) (2019,58) (2020,51)
        (2021,74) (2022,66) (2023,75) (2024,115) (2025,101) (2026,119)};
    \end{axis}
  \end{tikzpicture}
  \caption{Estudos elegíveis após a triagem, por ano de publicação ($n = \nElegiveis$; 2026 até à data das pesquisas). Destes, \nIncluidos{} foram analisados em texto integral.}
  \label{fig:por-ano}
\end{figure}

\paragraph{Tipo de publicação (QE2).} Dos \nIncluidos{} estudos analisados, dez ($52{,}6\,\%$) foram publicados
em revistas, seis ($31{,}6\,\%$) em conferências e três ($15{,}8\,\%$) em
\emph{workshops}.

% ---------------------------------------------------------------------------
\subsection{Taxonomia das soluções (QF2)}
\label{sec:taxonomia}

Os estudos foram agrupados em oito categorias segundo o ponto de atuação e o
objetivo do seu contributo principal (Figura~\ref{fig:taxonomia}).

\begin{figure}[ht]
  \centering
  \begin{tikzpicture}[
      raiz/.style={draw=iscte, fill=iscte, text=white, rounded corners=2pt,
                   font=\small\bfseries, minimum height=0.8cm, text width=4cm, align=center},
      ramo/.style={draw=iscte, fill=fundotab, rounded corners=2pt, font=\scriptsize\bfseries,
                   text width=3.3cm, minimum height=0.7cm, align=center},
      folha/.style={draw=cinza, rounded corners=2pt, font=\scriptsize, text width=3.3cm,
                    align=left, minimum height=0.6cm},
      ligacao/.style={draw=iscte, thick}]
    \node[raiz] (r) {Segurança em MQTT ($n=\nIncluidos$)};
    \node[ramo, below left=1.0cm and 2.0cm of r] (a) {Compreender ameaças};
    \node[ramo, below=1.0cm of r] (b) {Proteger ligação e \emph{broker}};
    \node[ramo, below right=1.0cm and 2.0cm of r] (c) {Proteger a mensagem e os dados};
    \node[folha, below=0.35cm of a] (c1) {C1 Análise de vulnerabilidades e ataques (3)};
    \node[folha, below=0.2cm of c1] (c7) {C7 Deteção de intrusões (2)};
    \node[folha, below=0.35cm of b] (c2) {C2 Autenticação (1)};
    \node[folha, below=0.2cm of c2] (c3) {C3 Autorização e controlo de acesso (3)};
    \node[folha, below=0.2cm of c3] (c6) {C6 Proteção do \emph{broker} e do transporte (3)};
    \node[folha, below=0.35cm of c] (c5) {C5 Proteção leve fim-a-fim da mensagem (5)};
    \node[folha, below=0.2cm of c5] (c4) {C4 Cifra baseada em atributos (1)};
    \node[folha, below=0.2cm of c4] (c8) {C8 Privacidade e anonimato (1)};
    \draw[ligacao] (r.south) -- ++(0,-0.4) -| (a.north);
    \draw[ligacao] (r.south) -- (b.north);
    \draw[ligacao] (r.south) -- ++(0,-0.4) -| (c.north);
    \foreach \p/\f in {a/c1, b/c2, c/c5} \draw[ligacao] (\p) -- (\f);
    \draw[ligacao] (c1) -- (c7); \draw[ligacao] (c2) -- (c3); \draw[ligacao] (c3) -- (c6);
    \draw[ligacao] (c5) -- (c4); \draw[ligacao] (c4) -- (c8);
  \end{tikzpicture}
  \caption{Taxonomia dos estudos incluídos, com o número de estudos por categoria.}
  \label{fig:taxonomia}
\end{figure}

\subsubsection{C1 --- Análise de vulnerabilidades e ataques}

Andy et al.~\cite{andy2017attack} mostram que, por omissão, o MQTT não cifra os
dados e demonstram ataques à privacidade, à autenticação e à integridade. Hernández
Ramos et al.~\cite{hernandez2018fuzzing} propõem uma técnica de \emph{fuzzing}
baseada em \emph{templates} que encontrou falhas novas em implementações reais. Vaccari et
al.~\cite{vaccari2020slowite} propõem o SlowITe, um ataque de negação de serviço de
baixa taxa, eficaz contra serviços reais com e sem cifra.

\subsubsection{C2 --- Autenticação}

Javadi et al.~\cite{javadi2026cmap} propõem o CMAP, autenticação mútua sem
certificados, sobre curvas elípticas, entre cada dispositivo e o \emph{broker}. A
segurança é demonstrada no modelo do oráculo aleatório e o protocolo é avaliado em
Raspberry Pi (338~ms de latência de autenticação contra 713~ms com TLS) e em
simulação NS-3. A resistência à repetição que reportam diz respeito às mensagens de
autenticação, protegidas por \emph{nonces} frescos, e não às mensagens publicadas.

\subsubsection{C3 --- Autorização e controlo de acesso}

Niruntasukrat et al.~\cite{niruntasukrat2016authorization} adaptam o OAuth~1.0a a
nós restritos e mostram atrasos de autorização mínimos numa plataforma real.
Rizzardi et al.~\cite{rizzardi2016aups} estendem o MQTT com gestão de chaves e
aplicação de políticas (AUPS, em código aberto). Calabretta et
al.~\cite{calabretta2018mqttauth} propõem o MQTT-Auth: o AugPAKE garante a
confidencialidade e dois \emph{tokens} autenticam e autorizam o acesso a cada
tópico.

\subsubsection{C4 --- Cifra baseada em atributos}

Singh et al.~\cite{singh2015secure} propõem o SMQTT e o SMQTT-SN, que acrescentam
ao MQTT cifra baseada em atributos com política na chave ou no criptograma
(KP/CP-ABE) sobre criptografia de curvas elípticas (ECC) leve, avaliada em simulação: só os subscritores com
atributos compatíveis decifram. A ABE resolve a confidencialidade de grupo, mas com
um custo computacional elevado para dispositivos restritos.

\subsubsection{C5 --- Proteção leve fim-a-fim da mensagem}

É a categoria mais numerosa. Dinculeană e Cheng~\cite{dinculeana2019vulnerabilities}
enviam valores sob a forma de HMAC (\emph{Value-to-HMAC}), com melhor desempenho do
que uma cifra simétrica; os próprios autores remetem para outros mecanismos a
autenticação dos clientes, a distribuição de chaves e a proteção contra
repetição~\cite[§6]{dinculeana2019vulnerabilities}. Malina et
al.~\cite{malina2019secure} propõem três níveis de segurança, do simples controlo de
adulteração à autenticação mútua de todas as partes. Todos assinam com Schnorr uma
entrada que inclui o nível, o tópico, a identidade do publicador, os dados e um
carimbo temporal~\cite[§4.5, Eq.~6]{malina2019secure}, o que dá não repúdio e liga a
mensagem ao tópico. Nos níveis 2 e 3 o carimbo serve de proteção contra
repetição~\cite[§5.2]{malina2019secure}, mas o \emph{broker} decifra o
conteúdo~\cite[§4.6]{malina2019secure}. Sadio et al.~\cite{sadio2019lightweight} usam
ChaCha20-Poly1305 fim-a-fim, com um protótipo em Arduino UNO. Dahlmanns et
al.~\cite{dahlmanns2021entrust} propõem o ENTRUST: um servidor de chaves distribui,
segundo uma ACL, uma chave por tópico aos clientes autorizados, sem alterar o
protocolo nem confiar no \emph{broker}~\cite[§5.2--5.3]{dahlmanns2021entrust}. Cada
mensagem leva um contador por publicador, que permite detetar duplicados e
perdas~\cite[§5.3]{dahlmanns2021entrust}. Van Glabbeek et
al.~\cite{vanglabbeek2022flexible} propõem três opções (autenticação e não repúdio;
mais confidencialidade; mais anonimato), todas com assinatura Schnorr e carimbo
temporal, com um \emph{broker} honesto mas curioso e uma autoridade de
certificação (CA) totalmente confiável~\cite[§3.1, §4.3]{vanglabbeek2022flexible}. Face ao MQTT sobre TLS, reduzem
em 80\,\% o \emph{overhead} de comunicação e em 40\,\% o computacional no
estabelecimento da ligação.

\subsubsection{C6 --- Proteção do \emph{broker} e do transporte}

Segarra et al.~\cite{segarra2020mqttz} executam o \emph{broker} num ambiente de
execução confiável (ARM TrustZone), avaliado com cargas reais de uma aplicação
médica. Spielvogel et al.~\cite{spielvogel2021tls} observam que um \emph{broker}
reencaminha mensagens independentemente da proteção de cada salto e propõem que
publicadores e subscritores possam exigir um nível de segurança mínimo. Prantl et
al.~\cite{prantl2021performance} medem o impacto do TLS num microcontrolador e
concluem que depende sobretudo das condições da rede e da frequência de
restabelecimento de ligações.

\subsubsection{C7 --- Deteção de intrusões}

Vaccari et al.~\cite{vaccari2020mqttset} publicam o MQTTset, um conjunto de dados
de tráfego MQTT legítimo e de ataques para treinar detetores. Alotaibi et
al.~\cite{alotaibi2024secure} detetam anomalias em tempo real com um enquadramento
distribuído (H2O), com os melhores resultados obtidos por XGBoost.

\subsubsection{C8 --- Privacidade e anonimato}

Buccafurri et al.~\cite{buccafurri2023mqtta} propõem o MQTT-A, uma rede de
\emph{brokers} intermédios que impede o \emph{broker} final de identificar a origem
e o destino das mensagens, sem alterar a infraestrutura MQTT.

% ---------------------------------------------------------------------------
\subsection{Comparação dos mecanismos (QF3, QF4)}
\label{sec:comparacao}

A Tabela~\ref{tab:comparacao} compara os 16 estudos que propõem mecanismos
(C2 a C8); os três estudos de análise (C1) ficam de fora. Cada célula foi
classificada a partir do texto integral, e a secção ou página que a sustenta está
registada no ficheiro de rastreabilidade do repositório~\cite{mqttsecure2026repo}.
\enquote{Não reportado} (\Nr) não equivale a \enquote{ausente por desenho}
(\Nao).

\begin{table}[p]
  \centering
  \caption{Comparação dos estudos que propõem mecanismos. Propriedades: C
    confidencialidade, I integridade, Au autenticação, Az autorização, V validade
    temporal, R anti-repetição, NR não repúdio, D disponibilidade/deteção, P privacidade.
    BnC: funciona com \emph{broker} não confiável. Aval.: E experimental, F formal,
    S simulação, A analítica. Restr.: avaliado ou desenhado para dispositivos
    restritos. \Sim{} garantido; \Parc{} parcialmente; \Nao{} não, por desenho;
    \Nr{} não reportado.}
  \label{tab:comparacao}
  \scriptsize
  \setlength{\tabcolsep}{2.6pt}
  \renewcommand{\arraystretch}{1.25}
  \begin{tabular}{@{}lL{2.55cm}L{3.55cm}ccccccccc|ccc@{}}
    \toprule
    \textbf{ID} & \textbf{Estudo} & \textbf{Mecanismo principal}
      & \textbf{C} & \textbf{I} & \textbf{Au} & \textbf{Az} & \textbf{V} & \textbf{R} & \textbf{NR}
      & \textbf{D} & \textbf{P} & \textbf{BnC} & \textbf{Aval.} & \textbf{Restr.} \\
    \midrule
    \multicolumn{15}{@{}l}{\textit{C2 --- Autenticação}} \\
    \E{19} & Javadi et al.~\cite{javadi2026cmap}            & Autenticação mútua sem certificados (ECC)
      & \Parc & \Parc & \Sim & \Nr & \Nr & \Nr & \Nr & \Nr & \Nr & \Nao & F, E, S & \Sim \\
    \multicolumn{15}{@{}l}{\textit{C3 --- Autorização e controlo de acesso}} \\
    \E{02} & Niruntasukrat et al.~\cite{niruntasukrat2016authorization} & OAuth 1.0a adaptado
      & \Nr & \Nr & \Parc & \Sim & \Nr & \Nr & \Nr & \Nr & \Nr & \Nr & E & \Sim \\
    \E{03} & Rizzardi et al.~\cite{rizzardi2016aups}        & Gestão de chaves e políticas (AUPS)
      & \Parc & \Nr & \Nr & \Sim & \Nr & \Nr & \Nr & \Nr & \Nr & \Nr & A, E & \Sim \\
    \E{05} & Calabretta et al.~\cite{calabretta2018mqttauth} & AugPAKE e \emph{tokens} por tópico
      & \Sim & \Nr & \Sim & \Sim & \Nr & \Nr & \Nr & \Nr & \Nr & \Nao & A, E & \Sim \\
    \multicolumn{15}{@{}l}{\textit{C4 --- Cifra baseada em atributos}} \\
    \E{01} & Singh et al.~\cite{singh2015secure}            & KP/CP-ABE sobre ECC
      & \Sim & \Nr & \Nr & \Sim & \Nr & \Nr & \Nr & \Nr & \Nr & \Nao & S & \Sim \\
    \multicolumn{15}{@{}l}{\textit{C5 --- Proteção leve fim-a-fim da mensagem}} \\
    \E{07} & Dinculeană e Cheng~\cite{dinculeana2019vulnerabilities} & Mapeamento \emph{Value-to-HMAC}
      & \Sim & \Sim & \Nr & \Nr & \Nr & \Nao & \Nr & \Nr & \Nr & \Nr & E & \Sim \\
    \E{08} & Malina et al.~\cite{malina2019secure}          & Três níveis de segurança (Schnorr)
      & \Parc & \Sim & \Sim & \Parc & \Parc & \Parc & \Sim & \Nr & \Sim & \Nao & A & \Sim \\
    \E{09} & Sadio et al.~\cite{sadio2019lightweight}       & AEAD ChaCha20-Poly1305
      & \Sim & \Sim & \Nr & \Nr & \Nr & \Nr & \Nr & \Nr & \Nr & \Sim & E & \Sim \\
    \E{13} & Dahlmanns et al.~\cite{dahlmanns2021entrust}   & Servidor de chaves por tópico (ENTRUST)
      & \Sim & \Sim & \Sim & \Sim & \Nr & \Sim & \Nr & \Parc & \Nr & \Sim & E & \Sim \\
    \E{16} & Van Glabbeek et al.~\cite{vanglabbeek2022flexible} & Três opções com assinatura Schnorr
      & \Parc & \Sim & \Sim & \Nr & \Parc & \Parc & \Sim & \Nr & \Sim & \Parc & E & \Sim \\
    \multicolumn{15}{@{}l}{\textit{C6 --- Proteção do broker e do transporte}} \\
    \E{12} & Segarra et al.~\cite{segarra2020mqttz}         & \emph{Broker} em TEE (TrustZone)
      & \Sim & \Nr & \Sim & \Parc & \Nr & \Nr & \Nr & \Nr & \Nr & \Parc & E & \Nr \\
    \E{14} & Spielvogel et al.~\cite{spielvogel2021tls}     & Nível de segurança exigido por salto
      & \Parc & \Nr & \Nr & \Nr & \Nr & \Nr & \Nr & \Nr & \Nr & \Nao & E & \Nr \\
    \E{15} & Prantl et al.~\cite{prantl2021performance}     & Avaliação do MQTT sobre TLS
      & \Sim & \Sim & \Nr & \Nr & \Nr & \Nr & \Nr & \Nr & \Nr & \Nao & E & \Sim \\
    \multicolumn{15}{@{}l}{\textit{C7 --- Deteção de intrusões}} \\
    \E{11} & Vaccari et al.~\cite{vaccari2020mqttset}       & Conjunto de dados para IDS
      & \Nr & \Nr & \Nr & \Nr & \Nr & \Nr & \Nr & \Parc & \Nr & \Nr & E & \Nr \\
    \E{18} & Alotaibi et al.~\cite{alotaibi2024secure}      & IDS distribuído (H2O, XGBoost)
      & \Nr & \Nr & \Nr & \Nr & \Nr & \Nr & \Nr & \Parc & \Nr & \Nr & E & \Nr \\
    \multicolumn{15}{@{}l}{\textit{C8 --- Privacidade e anonimato}} \\
    \E{17} & Buccafurri et al.~\cite{buccafurri2023mqtta}   & Rede de \emph{bridges} anónima
      & \Nr & \Nr & \Nr & \Nr & \Nr & \Nr & \Nr & \Nr & \Sim & \Nr & E & \Nr \\
    \bottomrule
  \end{tabular}
\end{table}

\paragraph{Cobertura das propriedades (QE3).} A Figura~\ref{fig:propriedades}
contabiliza, para os 16 mecanismos, quantos garantem cada propriedade (total ou
parcialmente). A confidencialidade é a propriedade mais tratada (12 estudos),
seguida da autorização (7), da autenticação (7) e da integridade (7). Só quatro
mecanismos funcionam, total ou parcialmente, com um \emph{broker} não confiável. A
validade temporal, a proteção contra repetição e o não repúdio são raros, mas não
estão ausentes: o ENTRUST deteta duplicados com um contador por
publicador~\cite{dahlmanns2021entrust}, e Malina et al.\ e Van Glabbeek et al.\
assinam um carimbo temporal, o que dá não repúdio e uma proteção parcial contra
repetição~\cite{malina2019secure, vanglabbeek2022flexible}. Nenhum dos três combina
um limite de idade com a deteção de duplicados dentro desse limite.

\begin{figure}[ht]
  \centering
  \begin{tikzpicture}
    \begin{axis}[
        xbar stacked, bar width=9pt, width=0.8\textwidth, height=6.4cm,
        xmin=0, xmax=18, xtick={0,2,...,18},
        symbolic y coords={NR,R,V,D,P,BnC,I,Az,Au,C},
        ytick=data, y tick label style={font=\small},
        xlabel={N.º de mecanismos (de 16)}, xlabel style={font=\small},
        legend style={font=\scriptsize, at={(0.98,0.04)}, anchor=south east},
        axis lines*=left, xmajorgrids, grid style={cinza!25}]
      \addplot[fill=acento!80, draw=iscte] coordinates {
        (7,C) (6,Au) (5,Az) (6,I) (2,BnC) (3,P) (0,D) (0,V) (1,R) (2,NR)};
      \addplot[fill=parcial!60, draw=parcial] coordinates {
        (5,C) (1,Au) (2,Az) (1,I) (2,BnC) (0,P) (3,D) (2,V) (2,R) (0,NR)};
      \legend{Garantida, Parcialmente}
    \end{axis}
  \end{tikzpicture}
  \caption{Número de mecanismos que garantem cada propriedade de segurança
    (siglas como na Tabela~\ref{tab:comparacao}).}
  \label{fig:propriedades}
\end{figure}

\paragraph{Avaliação (QF4).} Catorze dos 16 mecanismos ($87{,}5\,\%$) reportam
avaliação experimental; um é avaliado só em simulação~\cite{singh2015secure}, um só
analiticamente, por contagem de operações~\cite{malina2019secure}, e só um apresenta
uma prova formal de segurança~\cite{javadi2026cmap}. Onze são desenhados ou
avaliados para dispositivos restritos, do Arduino UNO~\cite{sadio2019lightweight} e
do Cortex-M3~\cite{dahlmanns2021entrust} ao Raspberry Pi~\cite{javadi2026cmap}. Nenhum dos estudos analisados mede, no mesmo ensaio, o custo de cada mecanismo e os
ataques que ele impede.

% ---------------------------------------------------------------------------
\subsection{Respostas às perguntas de investigação}
\label{sec:respostas}

\paragraph{QG e QF2.} Compreender as ameaças, proteger a ligação e o
\emph{broker}, e proteger a mensagem, sendo esta última a frente mais numerosa
(Secção~\ref{sec:taxonomia}).

\paragraph{QF1.} Escuta e injeção quando o MQTT é usado sem cifra, o que acontece
por omissão~\cite{andy2017attack}, negação de serviço que a cifra não
evita~\cite{vaccari2020slowite}, erros de implementação~\cite{hernandez2018fuzzing}
e o \emph{broker} como ponto de exposição de todo o
conteúdo~\cite{segarra2020mqttz, dahlmanns2021entrust}.

\paragraph{QF3 e QE3.} Ver a Figura~\ref{fig:propriedades}: a confidencialidade
domina, só quatro mecanismos dispensam a confiança no \emph{broker}, e a validade
temporal, a anti-repetição e o não repúdio aparecem isolados em três estudos.

\paragraph{QF4.} Sobretudo experimental (14 de 16); uma única prova formal.

\paragraph{QF5.} O custo da segurança em dispositivos a bateria~\cite{hintaw2023mqtt},
a inadequação do TLS a nós restritos~\cite{sadio2019lightweight, javadi2026cmap},
a exposição do conteúdo no \emph{broker}~\cite{segarra2020mqttz} e a impossibilidade
de um cliente impor a proteção dos restantes saltos~\cite{spielvogel2021tls}.

\section{Proposta: arquitetura de segurança em camadas}
\label{sec:proposta}

Esta secção propõe uma arquitetura que responde às lacunas identificadas na
Secção~\ref{sec:resultados}. A proposta é apresentada ao nível do desenho e da sua
formalização; a avaliação experimental (custo de cada camada e ataques que cada uma
impede) é objeto da segunda parte do projeto. Uma implementação de referência em
C++23, com documentação técnica completa, está disponível no repositório do
projeto~\cite{mqttsecure2026repo}.

% ---------------------------------------------------------------------------
\subsection{Princípios de desenho}

\begin{enumerate}[label=\textbf{D\arabic*.}, leftmargin=*]
  \item \textbf{Defesa em profundidade}~(Eq.~\ref{eq:profundidade}): proteção no
    transporte, na mensagem e na disponibilidade, de forma independente.
  \item \textbf{O \emph{broker} não é de confiança para o conteúdo}: a
    confidencialidade, a integridade, a autenticidade e a frescura das mensagens
    não podem depender dele.
  \item \textbf{Camadas independentes e comutáveis}: cada mecanismo é uma camada
    que pode ser ativada ou desativada sem alterar as restantes, o que permite medir
    isoladamente o custo e o benefício de cada um.
  \item \textbf{Verificações baratas primeiro, estado só depois da autenticação}:
    rejeitar cedo o que é barato rejeitar e nunca deixar entrada não autenticada
    alterar estado.
  \item \textbf{Estado limitado}: toda a estrutura que um atacante consiga fazer
    crescer tem um limite superior.
\end{enumerate}

% ---------------------------------------------------------------------------
\subsection{Arquitetura em camadas}
\label{sec:arquitetura}

A comunicação de cada extremo (sensor, \emph{gateway}) passa por uma sequência
ordenada de camadas entre a aplicação e o transporte MQTT
(Figura~\ref{fig:camadas}). Seja $\mathcal{M}$ o conjunto das mensagens e
$\bot$ a rejeição. Cada camada $\ell_i$ define duas funções
$\ell_i^{\uparrow}, \ell_i^{\downarrow} : \mathcal{M} \to \mathcal{M} \cup \{\bot\}$,
com $\ell(\bot) = \bot$. Para camadas ordenadas do lado da aplicação ($\ell_1$)
para o lado do transporte ($\ell_n$):
\begin{align}
  \operatorname{Saída}(m)  &= \ell_n^{\downarrow} \circ \cdots \circ \ell_1^{\downarrow}(m), \label{eq:saida}\\
  \operatorname{Entrada}(m) &= \ell_1^{\uparrow} \circ \cdots \circ \ell_n^{\uparrow}(m). \label{eq:entrada}
\end{align}
A primeira rejeição interrompe o processamento: uma mensagem recebida que alguma
camada rejeite nunca chega à aplicação. A \emph{ordem} das camadas tem significado:
com a persistência acima da segurança guarda-se texto em claro; abaixo dela,
guardam-se envelopes (cifra em repouso); a limitação de tráfego fica junto ao
transporte para descartar inundações antes de qualquer operação criptográfica.

\begin{figure}[ht]
  \centering
  \begin{tikzpicture}[
      camada/.style={draw=iscte, fill=fundotab, rounded corners=2pt, minimum width=3.2cm,
                     minimum height=0.75cm, font=\small, align=center},
      ext/.style={draw=cinza, fill=white, rounded corners=2pt, minimum width=3.2cm,
                  minimum height=0.75cm, font=\small, align=center},
      seta/.style={-{Stealth[length=2mm]}, thick, iscte},
      nota/.style={font=\scriptsize, text=cinza, align=left, text width=7.4cm}]
    \node[ext] (app) {Aplicação};
    \node[camada, below=0.35cm of app] (pers) {Persistência (opcional)};
    \node[camada, below=0.35cm of pers] (sec) {\textbf{Segurança da mensagem}};
    \node[camada, below=0.6cm of sec] (lim) {Limitação de tráfego};
    \node[ext, below=0.35cm of lim] (mqtt) {Transporte MQTT\\\scriptsize TLS / mTLS};
    \node[ext, below=0.6cm of mqtt, fill=iscte!10] (brk) {\emph{Broker}\\\scriptsize autenticação, ACL, limites};
    \foreach \a/\b in {app/pers, pers/sec, sec/lim, lim/mqtt}
      \draw[seta] ([xshift=-0.6cm]\a.south) -- ([xshift=-0.6cm]\b.north);
    \foreach \a/\b in {pers/app, sec/pers, lim/sec, mqtt/lim}
      \draw[seta, acento] ([xshift=0.6cm]\a.north) -- ([xshift=0.6cm]\b.south);
    \draw[seta, <->] (mqtt) -- (brk);
    \node[font=\scriptsize, text=iscte, left=0.1cm of pers.west, anchor=east] {saída $\downarrow$};
    \node[font=\scriptsize, text=acento, right=0.1cm of pers.east, anchor=west] {$\uparrow$ entrada};
    \node[nota, right=0.9cm of sec] (n1) {Envelope fim-a-fim: AES-256-GCM, GMAC ou
      Ed25519; tópico nos dados associados; frescura; janela anti-repetição;
      política de modos; restrição chave--tópico.};
    \node[nota, right=0.9cm of lim] (n2) {\emph{Token bucket} por tópico,
      antes de qualquer operação criptográfica; estado limitado.};
    \node[nota, right=0.9cm of brk] (n3) {Camada 1: TLS 1.3, mTLS com identidade no
      certificado, palavras-passe, ACL por utilizador, identificador de cliente
      ligado ao utilizador, limites de pacote e de ligações.};
    \draw[cinza, densely dotted] (sec.east) -- (n1.west);
    \draw[cinza, densely dotted] (lim.east) -- (n2.west);
    \draw[cinza, densely dotted] (brk.east) -- (n3.west);
  \end{tikzpicture}
  \caption{Arquitetura proposta: camadas entre a aplicação e o transporte, cada uma
    comutável de forma independente. A saída percorre as camadas de cima para
    baixo e a entrada de baixo para cima.}
  \label{fig:camadas}
\end{figure}

% ---------------------------------------------------------------------------
\subsection{Camada de transporte e \emph{broker}}

A primeira camada usa o que os \emph{brokers} já oferecem
(Tabela~\ref{tab:seguranca-nativa}), configurado de forma restritiva: só a porta
TLS aberta (TLS~1.3); autenticação mútua por certificado, com o nome comum do
certificado como identidade MQTT; palavras-passe; ACL de negação por omissão com
regras parametrizadas pelo utilizador autenticado (por exemplo,
\texttt{restricted/sensors/\%u/\#} só para escrita); identificador de cliente
igual ao nome de utilizador, o que impede um cliente de tomar a sessão de outro;
desativação de \texttt{\$SYS}; e limites de tamanho de pacote e de número de
ligações. Esta camada protege a rede e restringe quem fala com o \emph{broker},
mas, por si só, deixa intactos os problemas P1 a P3.

% ---------------------------------------------------------------------------
\subsection{Envelope fim-a-fim}
\label{sec:envelope}

A camada de segurança substitui o conteúdo de cada mensagem por um
\emph{envelope} (Figura~\ref{fig:envelope}), com cabeçalho fixo de 25 octetos.

\begin{figure}[ht]
  \centering
  \begin{tikzpicture}[
      campo/.style={draw=iscte, minimum height=0.85cm, font=\scriptsize, align=center,
                    fill=fundotab},
      corpo/.style={draw=iscte, minimum height=0.85cm, font=\scriptsize, align=center,
                    fill=acento!20}]
    \node[campo, minimum width=0.9cm] (m)  {\texttt{"MS"}\\2};
    \node[campo, minimum width=0.8cm, right=0pt of m] (v) {ver.\\1};
    \node[campo, minimum width=0.9cm, right=0pt of v] (md) {modo\\1};
    \node[campo, minimum width=1.4cm, right=0pt of md] (k) {id. chave\\4};
    \node[campo, minimum width=1.8cm, right=0pt of k] (s) {sequência\\8};
    \node[campo, minimum width=1.8cm, right=0pt of s] (t) {instante (ms)\\8};
    \node[campo, minimum width=0.7cm, right=0pt of t] (n) {$|N|$\\1};
    \node[campo, minimum width=1.3cm, right=0pt of n] (nn) {\emph{nonce}\\$|N|$};
    \node[corpo, minimum width=3.6cm, right=0pt of nn] (b) {corpo\\variável};
    \draw[decorate, decoration={brace, amplitude=5pt}, thick, bom]
      ([yshift=2pt]m.north west) -- ([yshift=2pt]nn.north east)
      node[midway, above=6pt, font=\scriptsize, text=bom]
      {autenticado como dados associados, juntamente com o tópico MQTT};
    \draw[decorate, decoration={brace, amplitude=5pt, mirror}, thick, iscte]
      ([yshift=-2pt]m.south west) -- ([yshift=-2pt]n.south east)
      node[midway, below=6pt, font=\scriptsize, text=iscte] {cabeçalho fixo: 25 octetos};
  \end{tikzpicture}
  \caption{Formato do envelope (inteiros em \emph{big-endian}); o número sob cada
    campo é o seu tamanho em octetos.}
  \label{fig:envelope}
\end{figure}

\paragraph{Dados associados e ligação ao tópico.} Sendo $\mathit{cab}$ o cabeçalho,
$N$ o \emph{nonce} e $\tau$ o tópico MQTT, os dados associados são
\begin{equation}
  A = \mathit{cab} \,\|\, N \,\|\, \operatorname{u16}(|\tau|) \,\|\, \tau .
  \label{eq:aad}
\end{equation}
Como o tópico entra na autenticação mas não no envelope, uma mensagem válida
republicada noutro tópico falha a verificação sem qualquer custo de transmissão; e
como todo o cabeçalho (incluindo o modo, o identificador da chave, a sequência e o
instante) é autenticado, nenhum desses campos pode ser alterado.

\paragraph{Modos de proteção.} O remetente escolhe um de três modos, indicado no
cabeçalho:
\begin{align}
  \textbf{cifra:}\quad & \mathit{corpo} = C \,\|\, T, & (C, T) &= \operatorname{AES\text{-}GCM.Enc}(K, N, P, A); \label{eq:modo-cifra}\\
  \textbf{autenticação:}\quad & \mathit{corpo} = P \,\|\, T, & T &= \operatorname{AES\text{-}GCM.Enc}(K, N, \varnothing, A \,\|\, P); \label{eq:modo-auth}\\
  \textbf{assinatura:}\quad & \mathit{corpo} = P \,\|\, \sigma, & \sigma &= \operatorname{Ed25519.Sign}(\mathit{sk}, A \,\|\, P),\ N = \varnothing. \label{eq:modo-sign}
\end{align}
O modo de cifra garante confidencialidade e integridade; o de autenticação (GMAC)
garante integridade com o conteúdo legível; o de assinatura garante integridade e
\emph{prova de origem}: o recetor guarda apenas chaves públicas e, por isso, consegue
verificar, mas não forjar, mensagens dos dispositivos. O custo fixo por mensagem é
\begin{equation}
  \omega(\text{modo}) = 25 + |N| + |T| =
  \begin{cases}
    25 + 12 + 16 = 53 \text{ octetos} & \text{(cifra, autenticação)},\\
    25 + 0 + 64 = 89 \text{ octetos} & \text{(assinatura)}.
  \end{cases}
  \label{eq:overhead}
\end{equation}

\paragraph{Política do recetor contra a redução do nível de proteção.} Cada recetor
mantém o conjunto $\mathcal{A} \subseteq \{\text{cifra}, \text{autenticação},
\text{assinatura}\}$ dos modos que aceita e rejeita qualquer envelope cujo modo
$\mu \notin \mathcal{A}$. Como $\mu$ faz parte de $A$, um atacante não consegue
rotular uma mensagem cifrada como \enquote{apenas autenticada} sem invalidar a
etiqueta, e um detentor de chave não consegue baixar o nível de proteção abaixo do
que o recetor exige.

% ---------------------------------------------------------------------------
\subsection{Ordem de verificação no recetor}
\label{sec:ordem}

A Tabela~\ref{tab:ordem} define a sequência de verificações de um envelope
recebido. A ordem decorre do princípio D4.

\begin{table}[ht]
  \centering
  \caption{Sequência de verificações de um envelope recebido.}
  \label{tab:ordem}
  \small
  \begin{tabularx}{\textwidth}{@{}cL{3.5cm}YL{3.2cm}@{}}
    \toprule
    \textbf{\#} & \textbf{Verificação} & \textbf{Justificação da posição} & \textbf{Rejeição} \\
    \midrule
    1 & Formato do envelope & Descarta conteúdo não protegido (por exemplo, injeções em claro) & mensagem malformada \\
    2 & Modo $\mu \in \mathcal{A}$ & Barato; o modo está em $A$, pelo que mentir sobre ele só pode falhar & modo não aceite \\
    3 & Frescura, Eq.~\eqref{eq:frescura} & Barato; o instante está em $A$, pelo que alterá-lo invalida a etiqueta & mensagem expirada \\
    4 & Chave conhecida e verificação criptográfica & A partir daqui a mensagem é autêntica; a falha é sempre a mesma, sem revelar a causa (sem oráculo) & falha de autenticação \\
    5 & Restrição chave--tópico, Eq.~\eqref{eq:ligacao} & Só depois da autenticação: entrada não autenticada não pode sondar a política & permissão negada \\
    6 & Janela anti-repetição, Eq.~\eqref{eq:janela} & Só depois da autenticação (números forjados não avançam a janela) e da restrição (uma mensagem recusada não consome o seu número) & repetição detetada \\
    \bottomrule
  \end{tabularx}
\end{table}

% ---------------------------------------------------------------------------
\subsection{Frescura e anti-repetição}

A frescura usa a Eq.~\eqref{eq:frescura} com $\Delta_a = 60$~s e
$\Delta_s = 5$~s, e a anti-repetição a janela da Eq.~\eqref{eq:janela} com
$w = 64$, mantida por identificador de chave. Os dois mecanismos são
complementares: a frescura sozinha deixa passar todas as repetições com menos de
$\Delta_a$; a janela sozinha não protege contra a perda do seu estado (por exemplo,
reinício do recetor), situação em que a frescura limita a exposição a $\Delta_a$.
Para cumprir o princípio D5, o número de janelas é limitado e um identificador de
chave novo é recusado quando o limite é atingido. Os limites do GCM
(Eq.~\eqref{eq:aniversario}) e do espaço de sequências ($2^{64}$) impõem a rotação
das chaves.

% ---------------------------------------------------------------------------
\subsection{Restrição chave--tópico}
\label{sec:ligacao}

Autenticar uma mensagem prova apenas que foi produzida com a chave $k$. Para que um
dispositivo comprometido não publique, com a sua chave válida, nos tópicos dos
outros (problema P3), cada chave tem um dono $\operatorname{dono}(k)$ e o recetor
aplica regras de negação por omissão. Sendo $F_k$ as regras específicas da chave,
$F_\ast$ os padrões gerais e $\sigma_k$ a substituição
$\{\texttt{\%u} \mapsto \operatorname{dono}(k),\ \texttt{\%k} \mapsto k\}$:
\begin{equation}
  \operatorname{autorizada}(k, \tau) \iff
  \exists f \in F_k \cup \sigma_k(F_\ast) :\ \operatorname{corr}(f, \tau),
  \label{eq:ligacao}
\end{equation}
com $\operatorname{corr}$ definida na Eq.~\eqref{eq:corr}. Por exemplo, com o padrão
\texttt{+/sensors/\%u/\#}, a chave de \texttt{sensor-7} só pode publicar em
\texttt{open/sensors/sensor-7/\dots} e \texttt{restricted/sensors/sensor-7/\dots}.
O dono tem de ser um único nível de tópico sem carateres especiais, para que a
substituição não alargue o filtro.

% ---------------------------------------------------------------------------
\subsection{Gestão de chaves}

Cada dispositivo $d$ tem, por época $e$, uma chave AES, um par de chaves Ed25519 e o
identificador $\operatorname{id}(d, e) = e \cdot 10^{6} + d$ (o \emph{gateway} é
$d = 0$). Identificadores nunca são reutilizados: um identificador reutilizado teria
uma janela anti-repetição nova, na qual mensagens antigas voltariam a ser aceites. A
Tabela~\ref{tab:chaves} mostra que chaves cada papel detém.

\begin{table}[ht]
  \centering
  \caption{Distribuição de chaves por papel.}
  \label{tab:chaves}
  \small
  \begin{tabularx}{\textwidth}{@{}L{2.4cm}YYY@{}}
    \toprule
    \textbf{Papel} & \textbf{Chaves AES} & \textbf{Chaves privadas Ed25519} & \textbf{Chaves públicas Ed25519} \\
    \midrule
    \emph{Gateway}  & de todos os dispositivos e a sua & \textbf{apenas a sua} & de todos os dispositivos e a sua \\
    Dispositivo $d$ & a sua e a do \emph{gateway} (ordens) & a sua & a do \emph{gateway} \\
    Atacante        & --- & --- & --- (apenas o certificado da CA) \\
    \bottomrule
  \end{tabularx}
\end{table}

A \textbf{rotação} instala as chaves da época $e+1$ ao lado das da época $e$,
muda os remetentes para as novas e revoga as antigas após um período de graça. A
\textbf{revogação} de um dispositivo remove as suas chaves de mensagem, a sua
palavra-passe no \emph{broker} e coloca o seu certificado na CRL, atuando nas
duas camadas de defesa.

% ---------------------------------------------------------------------------
\subsection{Limitação de tráfego antes da criptografia}

Junto ao transporte, uma camada aplica a Eq.~\eqref{eq:token-bucket} por tópico,
rejeitando o excesso antes de qualquer verificação criptográfica. O número de baldes
é limitado; quando o limite é atingido, os baldes inativos (que voltariam a estar
cheios) são descartados e, se não houver nenhum, os tópicos novos partilham um único
balde, pelo que um atacante não consegue criar orçamento inventando tópicos.

% ---------------------------------------------------------------------------
\subsection{Mapeamento entre ameaças e mecanismos}

A Tabela~\ref{tab:ameacas-mecanismos} relaciona as ameaças da
Tabela~\ref{tab:stride} com os mecanismos da proposta.

\begin{table}[ht]
  \centering
  \caption{Ameaças e mecanismos da arquitetura proposta.}
  \label{tab:ameacas-mecanismos}
  \small
  \begin{tabularx}{\textwidth}{@{}L{5.0cm}YL{2.6cm}@{}}
    \toprule
    \textbf{Ameaça} & \textbf{Mecanismo} & \textbf{Camada} \\
    \midrule
    Escuta na rede               & TLS 1.3 & transporte \\
    Escuta pelo \emph{broker} ou por subscritores de \texttt{\#} & Modo de cifra; ACL & mensagem; \emph{broker} \\
    Adulteração do conteúdo ou do cabeçalho & Etiqueta GCM ou assinatura sobre $A \,\|\, P$ & mensagem \\
    Republicação noutro tópico   & Tópico em $A$, Eq.~\eqref{eq:aad} & mensagem \\
    Repetição                    & Janela anti-repetição & mensagem \\
    Mensagens antigas            & Frescura & mensagem \\
    Injeção de mensagens em claro & Formato do envelope e autenticação & mensagem \\
    Redução do nível de proteção & Modo em $A$ e política $\mathcal{A}$ & mensagem \\
    Usurpação de identidade no \emph{broker} & Palavras-passe, mTLS, identificador ligado ao utilizador & transporte \\
    Dispositivo comprometido a publicar por outros & Restrição chave--tópico, Eq.~\eqref{eq:ligacao}; ACL com \texttt{\%u} & mensagem; \emph{broker} \\
    Recetor a forjar dados de dispositivos & Modo de assinatura (recetor só com chaves públicas) & mensagem \\
    Chaves roubadas              & Rotação por épocas; revogação & gestão \\
    Inundação de um tópico       & \emph{Token bucket} antes da criptografia; limites no \emph{broker} & disponibilidade \\
    Esgotamento de memória       & Estado limitado (janelas, baldes, repositório) & todas \\
    Oráculo de verificação       & Falha uniforme; estado só depois da autenticação & mensagem \\
    \bottomrule
  \end{tabularx}
\end{table}

% ---------------------------------------------------------------------------
\subsection{Comparação com os trabalhos relacionados}
\label{sec:comparacao-proposta}

A Tabela~\ref{tab:proposta-vs} compara a proposta com os mecanismos fim-a-fim e de
autenticação mais próximos, acrescentando às propriedades da
Tabela~\ref{tab:comparacao} três mecanismos específicos desta proposta: ligação ao
tópico (LT), política contra a redução do nível de proteção (RN) e restrição
chave--tópico (CT).

\begin{table}[ht]
  \centering
  \caption{Posicionamento da proposta face aos trabalhos relacionados mais próximos.
    Siglas e símbolos como na Tabela~\ref{tab:comparacao}; LT ligação ao tópico;
    RN política contra a redução do nível de proteção; CT restrição chave--tópico;
    Com.: camadas comutáveis em tempo de execução.}
  \label{tab:proposta-vs}
  \scriptsize
  \setlength{\tabcolsep}{2.8pt}
  \renewcommand{\arraystretch}{1.2}
  \begin{tabular}{@{}lL{2.8cm}ccccccccc|cccc@{}}
    \toprule
    \textbf{ID} & \textbf{Trabalho}
      & \textbf{C} & \textbf{I} & \textbf{Au} & \textbf{Az} & \textbf{F/R} & \textbf{NR}
      & \textbf{D} & \textbf{P} & \textbf{BnC} & \textbf{LT} & \textbf{RN} & \textbf{CT} & \textbf{Com.} \\
    \midrule
    \E{02} & Singh et al.~\cite{singh2015secure}
      & \Sim & \Nr & \Nr & \Sim & \Nr & \Nr & \Nr & \Nr & \Parc & \Nr & \Nr & \Nr & \Nr \\
    \E{11} & Calabretta et al.~\cite{calabretta2018mqttauth}
      & \Sim & \Nr & \Sim & \Sim & \Nr & \Nr & \Nr & \Nr & \Nr & \Nr & \Nr & \Nr & \Nr \\
    \E{13} & Dinculeană e Cheng~\cite{dinculeana2019vulnerabilities}
      & \Sim & \Sim & \Nr & \Nr & \Nr & \Nr & \Nr & \Nr & \Sim & \Nr & \Nr & \Nr & \Nr \\
    \E{15} & Malina et al.~\cite{malina2019secure}
      & \Parc & \Sim & \Sim & \Nr & \Nr & \Nr & \Nr & \Sim & \Nr & \Nr & \Nr & \Nr & \Nr \\
    \E{16} & Sadio et al.~\cite{sadio2019lightweight}
      & \Sim & \Sim & \Nr & \Nr & \Nr & \Nr & \Nr & \Nr & \Sim & \Nr & \Nr & \Nr & \Nr \\
    \E{19} & Segarra et al.~\cite{segarra2020mqttz}
      & \Sim & \Nr & \Sim & \Nr & \Nr & \Nr & \Nr & \Nr & \Parc & \Nr & \Nr & \Nr & \Nr \\
    \E{20} & Dahlmanns et al.~\cite{dahlmanns2021entrust}
      & \Sim & \Sim & \Sim & \Nr & \Nr & \Nr & \Nr & \Nr & \Sim & \Nr & \Nr & \Nr & \Nr \\
    \E{26} & Javadi et al.~\cite{javadi2026cmap}
      & \Sim & \Sim & \Sim & \Nr & \Nr & \Nr & \Nr & \Nr & \Nr & \Nr & \Nr & \Nr & \Nr \\
    \midrule
    \rowcolor{fundotab}
    --- & \textbf{Proposta}
      & \Sim & \Sim & \Sim & \Sim & \Sim & \Sim & \Parc & \Nao & \Sim & \Sim & \Sim & \Sim & \Sim \\
    \bottomrule
  \end{tabular}
\end{table}

A proposta não supera os trabalhos relacionados em todas as dimensões: não oferece
anonimato~\cite{buccafurri2023mqtta}, não protege o próprio \emph{broker} com um
TEE~\cite{segarra2020mqttz}, não oferece confidencialidade de grupo por atributos
\cite{singh2015secure} e não tem verificação formal~\cite{javadi2026cmap}. O seu
contributo está na \emph{combinação}, numa só camada fim-a-fim, de propriedades
que a literatura analisada trata separadamente ou não reporta (ligação ao tópico,
frescura com anti-repetição, política contra a redução do nível de proteção,
restrição chave--tópico e não repúdio opcional) e na organização em camadas
comutáveis, que permite medir o custo e o benefício de cada mecanismo isoladamente.

% ---------------------------------------------------------------------------
\subsection{Limitações conhecidas}

\begin{itemize}[leftmargin=*]
  \item \textbf{Chaves pré-partilhadas}: a distribuição inicial das chaves assume um
    canal seguro; um acordo de chaves por dispositivo (por exemplo, ECDH) eliminaria
    os segredos pré-partilhados.
  \item \textbf{Chave de grupo nas ordens}: nos modos AES todos os atuadores detêm a
    chave do \emph{gateway}, pelo que um atuador comprometido pode forjar ordens para
    os outros; o modo de assinatura resolve-o, a um custo computacional superior.
  \item \textbf{Limitação por tópico}: o tópico é escolhido por quem publica, pelo
    que um atacante que inunde o tópico de uma vítima consome também o orçamento da
    vítima; limites por cliente pertencem ao \emph{broker}.
  \item \textbf{Estado da janela não persistente}: o reinício do recetor apaga as
    janelas anti-repetição, ficando a proteção limitada à frescura ($\Delta_a$).
  \item \textbf{Dados falsos mas autênticos}: um dispositivo comprometido continua a
    poder publicar dados falsos nos seus próprios tópicos; só a deteção e a
    revogação o travam.
\end{itemize}

% ---------------------------------------------------------------------------
\subsection{Plano de avaliação (segunda parte do projeto)}

A avaliação experimental medirá, para cada combinação de camadas: (i)~o custo, em
latência fim-a-fim, tempo de processador por mensagem, débito e octetos
acrescentados (Eq.~\ref{eq:overhead}); e (ii)~o benefício, como a fração de
ataques bem-sucedidos (escuta, injeção, repetição, mudança de tópico, usurpação de
identidade, inundação) com cada defesa ativa ou inativa. Esta medição conjunta de
custo e benefício responde diretamente à lacuna identificada na QF4.

\section{Discussão}
\label{sec:discussao}

\subsection{Lacunas do estado da arte}

A comparação das Secções~\ref{sec:resultados} e~\ref{sec:proposta} evidencia cinco
lacunas.

\begin{enumerate}[label=\textbf{L\arabic*.}, leftmargin=*]
  \item \textbf{A autenticidade raramente é acompanhada de frescura.} Nenhum dos
    mecanismos analisados reporta proteção contra repetição ou frescura como
    propriedade garantida. No MQTT, onde os duplicados legítimos do QoS~1 coexistem
    com repetições maliciosas, a ausência destes mecanismos permite reinjetar
    mensagens autênticas indefinidamente.
  \item \textbf{A mensagem não é ligada ao seu contexto.} A ligação criptográfica ao
    tópico, a proteção do modo contra a redução do nível de proteção e a restrição
    dos tópicos em que cada chave pode publicar não aparecem nos estudos
    analisados como mecanismos explícitos, apesar de corresponderem a ameaças
    concretas (mudança de tópico, \emph{downgrade}, dispositivo comprometido).
  \item \textbf{O \emph{broker} continua, em regra, a ser de confiança.} Só seis dos
    22 mecanismos dispensam, total ou parcialmente, a confiança no \emph{broker}; a
    maioria das soluções de autenticação e autorização reforça o \emph{broker} em
    vez de o tornar irrelevante para a segurança do conteúdo.
  \item \textbf{O não repúdio está ausente.} As soluções baseiam-se em segredos
    partilhados (palavras-passe, chaves simétricas, ABE), o que não impede quem
    detém o segredo --- incluindo o recetor --- de forjar mensagens em nome do
    remetente.
  \item \textbf{Custo e benefício não são medidos em conjunto.} As avaliações
    medem o custo do mecanismo, mas raramente o relacionam, no mesmo ensaio, com os
    ataques que o mecanismo impede; e os mecanismos não são comparáveis entre si
    por usarem plataformas, cargas e métricas diferentes.
\end{enumerate}

A arquitetura proposta responde diretamente a L1, L2, L3 e L4, e o seu desenho em
camadas comutáveis foi pensado para que a segunda parte do projeto responda a L5.

\subsection{Ameaças à validade}

Seguindo~\cite{kitchenham2007}, identificam-se as principais ameaças à validade
desta revisão e as medidas de mitigação adotadas (Tabela~\ref{tab:validade}).

\begin{table}[ht]
  \centering
  \caption{Ameaças à validade da revisão e mitigação.}
  \label{tab:validade}
  \small
  \begin{tabularx}{\textwidth}{@{}L{2.6cm}YY@{}}
    \toprule
    \textbf{Tipo} & \textbf{Ameaça} & \textbf{Mitigação} \\
    \midrule
    Construção & A \emph{string} pode não captar estudos que usem outra
      terminologia & Grupos de sinónimos; refinamento após a leitura das revisões;
      inclusão de MQTT-SN e de termos de deteção \\
    Seleção & Enviesamento na aplicação dos critérios & Dois revisores
      independentes; kappa de Cohen (Eq.~\ref{eq:kappa}); resolução das
      divergências por discussão \\
    Extração & Subjetividade na classificação das propriedades; um estudo pode
      garantir uma propriedade sem a descrever & Formulário fixo; símbolo
      \enquote{não reportado} distinto de \enquote{ausente por desenho};
      confirmação na leitura integral \\
    Externa & Bases de dados e língua limitam a cobertura; enviesamento de
      publicação & Cinco bases de dados; inglês e português; estudos de revista,
      conferência e \emph{workshop} \\
    Conclusão & Os números por propriedade dependem da granularidade das
      categorias & Taxonomia e critérios publicados; dados de extração no
      Anexo~\ref{anx:extracao} \\
    \bottomrule
  \end{tabularx}
\end{table}

A ameaça mais relevante é a da extração: as conclusões L1, L2 e L4 assentam em
propriedades \emph{não reportadas}, o que não equivale a propriedades
\emph{ausentes}. Por isso, a Tabela~\ref{tab:comparacao} distingue explicitamente
os dois casos e a discussão fala em propriedades \enquote{não reportadas}.

\section{Conclusões e trabalho futuro}
\label{sec:conclusoes}

\subsection{Conclusões}

Este trabalho apresentou uma revisão sistemática da literatura sobre mecanismos de
segurança para comunicações IoT baseadas em MQTT, conduzida segundo um protocolo
explícito --- perguntas de investigação, PICOC, \emph{strings} de pesquisa por base
de dados, critérios de inclusão, exclusão e qualidade, leitura em três passagens e
formulário de extração --- e relatada com o diagrama PRISMA. Foram analisados
\nIncluidos{} estudos primários (2014--2026) e \nSecundarios{} revisões
relacionadas.

As principais conclusões são as seguintes:

\begin{itemize}[leftmargin=*]
  \item O MQTT delega a segurança na implementação, e o TLS, por terminar no
    \emph{broker}, não protege as mensagens de um \emph{broker} comprometido nem de
    clientes com acesso legítimo, como mostram os estudos de análise de ameaças.
  \item A literatura organiza-se em oito categorias. A proteção leve fim-a-fim da
    mensagem é a mais numerosa, e a confidencialidade, a autenticação e a
    autorização são as propriedades mais tratadas.
  \item A frescura e a anti-repetição, a ligação da mensagem ao tópico, a política
    contra a redução do nível de proteção, a restrição chave--tópico e o não repúdio
    são pouco ou nada reportados, e a maioria das soluções continua a confiar no
    \emph{broker}.
  \item As avaliações medem o custo dos mecanismos, mas raramente o relacionam com
    os ataques efetivamente impedidos.
\end{itemize}

Com base nestas lacunas, propôs-se uma arquitetura de segurança em camadas
independentes e comutáveis, que combina defesa no transporte, um envelope fim-a-fim
formalmente especificado (AES-256-GCM, GMAC ou Ed25519, com o tópico e o cabeçalho
nos dados associados, frescura, janela anti-repetição, política de modos e
restrição chave--tópico) e limitação de tráfego antes de qualquer operação
criptográfica. A proposta não pretende substituir trabalhos com objetivos distintos,
como o anonimato ou a proteção do \emph{broker} por TEE, mas combinar numa só
camada propriedades que a literatura trata em separado.

\subsection{Trabalho futuro}

\begin{itemize}[leftmargin=*]
  \item \textbf{Avaliação experimental (segunda parte do projeto)}: medir, para cada
    combinação de camadas, o custo (latência, processamento, octetos) e o benefício
    (ataques impedidos) sobre um \emph{broker} real, respondendo à lacuna L5.
  \item \textbf{Verificação formal} das propriedades do envelope, à semelhança
    de~\cite{javadi2026cmap}.
  \item \textbf{Acordo de chaves por dispositivo} (por exemplo, ECDH autenticado),
    eliminando chaves pré-partilhadas, e persistência do estado anti-repetição.
  \item \textbf{Cifras para nós sem AES em \emph{hardware}}, como o
    ChaCha20-Poly1305~\cite{rfc8439, sadio2019lightweight}, e cifras resistentes à
    reutilização de \emph{nonces} para chaves de longa duração.
  \item \textbf{Integração com deteção de intrusões}~\cite{vaccari2020mqttset,
    alotaibi2024secure}, usando as rejeições de cada camada como sinal.
  \item \textbf{Atualização periódica da revisão}, reaplicando o protocolo da
    Secção~\ref{sec:metodologia}.
\end{itemize}


% ---------------------------------------------------- referências ---------
\clearpage
\bibliographystyle{IEEEtran}
\bibliography{referencias}

% -------------------------------------------------------- anexos ----------
\clearpage
\appendix
\renewcommand{\thesection}{\Alph{section}}
\section{\emph{Strings} de pesquisa por base de dados}
\label{anx:strings}

Todas as pesquisas foram restringidas a publicações de 2014 a 2026. Os grupos G1 a
G3 são os da Tabela~\ref{tab:grupos}; nas \emph{strings} secundárias acrescenta-se,
com \texttt{AND}, o grupo G4.

\newcommand{\str}[1]{\noindent\begin{minipage}{\linewidth}\footnotesize\ttfamily\raggedright #1\end{minipage}}

\subsection*{IEEE Xplore (\emph{Command Search}, campo \emph{All Metadata})}
\str{("All Metadata":"MQTT" OR "All Metadata":"MQTT-SN" OR "All Metadata":"MQ Telemetry Transport"
OR "All Metadata":"Message Queuing Telemetry Transport" OR "All Metadata":"Message Queue Telemetry Transport")
AND ("All Metadata":security OR "All Metadata":secure OR "All Metadata":authentication
OR "All Metadata":authorization OR "All Metadata":"access control" OR "All Metadata":encryption
OR "All Metadata":confidentiality OR "All Metadata":integrity OR "All Metadata":privacy
OR "All Metadata":attack* OR "All Metadata":vulnerabilit* OR "All Metadata":intrusion)
AND ("All Metadata":"Internet of Things" OR "All Metadata":IoT OR "All Metadata":IIoT
OR "All Metadata":"publish/subscribe" OR "All Metadata":"publish-subscribe"
OR "All Metadata":"machine-to-machine" OR "All Metadata":M2M)\\
Filtro: Year 2014--2026; Conferences, Journals.}

\subsection*{Scopus}
\str{TITLE-ABS-KEY ( ( "MQTT" OR "MQTT-SN" OR "MQ Telemetry Transport"
OR "Message Queuing Telemetry Transport" OR "Message Queue Telemetry Transport" )
AND ( "security" OR "secure" OR "authentication" OR "authorization" OR "access control"
OR "encryption" OR "confidentiality" OR "integrity" OR "privacy" OR attack* OR vulnerabilit*
OR "intrusion" )
AND ( "Internet of Things" OR "IoT" OR "IIoT" OR "publish/subscribe" OR "publish-subscribe"
OR "machine-to-machine" OR "M2M" ) )
AND PUBYEAR > 2013 AND PUBYEAR < 2027
AND ( LIMIT-TO ( DOCTYPE , "ar" ) OR LIMIT-TO ( DOCTYPE , "cp" ) )
AND ( LIMIT-TO ( LANGUAGE , "English" ) OR LIMIT-TO ( LANGUAGE , "Portuguese" ) )}

\subsection*{ACM Digital Library (pesquisa avançada, campos \emph{Title}, \emph{Abstract}, \emph{Keywords})}
\str{[[Title: "mqtt"] OR [Abstract: "mqtt"] OR [Keywords: "mqtt"]
OR [Abstract: "message queuing telemetry transport"] OR [Abstract: "mq telemetry transport"]]
AND [[Abstract: security] OR [Abstract: secure] OR [Abstract: authentication]
OR [Abstract: authorization] OR [Abstract: "access control"] OR [Abstract: encryption]
OR [Abstract: confidentiality] OR [Abstract: integrity] OR [Abstract: privacy]
OR [Abstract: attack*] OR [Abstract: vulnerabilit*] OR [Abstract: intrusion]]
AND [[Abstract: "internet of things"] OR [Abstract: iot] OR [Abstract: iiot]
OR [Abstract: "publish/subscribe"] OR [Abstract: "publish-subscribe"]
OR [Abstract: "machine-to-machine"] OR [Abstract: m2m]]\\
Filtro: Publication Date 2014--2026.}

\subsection*{Web of Science (Core Collection, campo \emph{Topic})}
\str{TS=( ("MQTT" OR "MQTT-SN" OR "MQ Telemetry Transport" OR "Message Queuing Telemetry Transport"
OR "Message Queue Telemetry Transport")
AND (security OR secure OR authentication OR authorization OR "access control" OR encryption
OR confidentiality OR integrity OR privacy OR attack* OR vulnerabilit* OR intrusion)
AND ("Internet of Things" OR IoT OR IIoT OR "publish/subscribe" OR "publish-subscribe"
OR "machine-to-machine" OR M2M) )\\
Filtro: PY=(2014-2026); Document Types: Article, Proceedings Paper.}

\subsection*{ScienceDirect (campo \emph{Title, abstract or author-specified keywords})}
O ScienceDirect não suporta o carácter universal \texttt{*} e limita o número de
operadores booleanos por campo; a \emph{string} foi, por isso, dividida em duas
consultas cujos resultados foram unidos:
\par\medskip
\str{(1) ("MQTT" OR "Message Queuing Telemetry Transport") AND (security OR authentication
OR authorization OR encryption) AND ("Internet of Things" OR IoT)\\[4pt]
(2) ("MQTT" OR "Message Queuing Telemetry Transport") AND (attack OR attacks OR vulnerability
OR intrusion OR privacy) AND ("Internet of Things" OR IoT)\\
Filtro: Years 2014--2026; Article types: Research articles.}

\subsection*{\emph{String} secundária (revisões), em todas as bases}
\str{<\emph{string} principal> AND ("systematic literature" OR "systematic review"
OR "literature review" OR "survey")}

\section{Extração de dados e referências dos estudos}
\label{anx:extracao}

A Tabela~\ref{tab:extracao} resume a extração de dados de cada estudo analisado e
identifica-o pela publicação, pelo DOI e pela versão do texto integral consultada;
as secções citadas referem-se a essa versão. A localização de cada célula das
tabelas~\ref{tab:comparacao} e~\ref{tab:proposta-vs} está no ficheiro
\texttt{report/rsl/rastreabilidade.md} do repositório.

{\footnotesize
\setlength{\tabcolsep}{3pt}
\begin{longtable}{@{}lL{2.2cm}cL{4.4cm}L{3.0cm}L{3.6cm}@{}}
\caption{Extração de dados, publicação e versão consultada dos estudos e revisões analisados.}\label{tab:extracao}\\
\toprule
\textbf{ID} & \textbf{Estudo} & \textbf{Cat.} & \textbf{Contributo principal} & \textbf{Avaliação reportada} & \textbf{Publicação; DOI; texto consultado} \\
\midrule
\endfirsthead
\toprule
\textbf{ID} & \textbf{Estudo} & \textbf{Cat.} & \textbf{Contributo principal} & \textbf{Avaliação reportada} & \textbf{Publicação; DOI; texto consultado} \\
\midrule
\endhead
\bottomrule
\endfoot
\E{01} & Singh et al.~\cite{singh2015secure} & C4 &
  SMQTT e SMQTT-SN: KP/CP-ABE sobre ECC leve; o \emph{broker} é o gerador de chaves (§III-B) & Três portáteis a simular dispositivos IoT (§VII-A) & CSNT 2015, pp.~746--751; \url{10.1109/CSNT.2015.16}; Editor (IEEE Xplore) \\
\E{02} & Niruntasukrat et al.~\cite{niruntasukrat2016authorization} & C3 &
  Autorização baseada em OAuth~1.0a adaptada a nós restritos sem interface de utilizador & Plataforma IoT real; atraso de autorização em Arduino Mega e outros (§IV) & IEEE ICC Workshops 2016, pp.~290--295; \url{10.1109/ICCW.2016.7503802}; Editor (IEEE Xplore) \\
\E{03} & Rizzardi et al.~\cite{rizzardi2016aups} & C3 &
  AUPS: gestão de chaves por tópico e aplicação de políticas sobre o MQTT; código aberto (Apache~2.0) & Robustez (§IV-A); \emph{overhead}, latência e computação em Raspberry Pi (§V) & Information Systems 62 (2016), pp.~29--41; \url{10.1016/j.is.2016.05.004}; Versão de autor (15~pp., paginação própria) \\
\E{04} & Andy et al.~\cite{andy2017attack} & C1 &
  Cenários de ataque ao MQTT sem cifra; razões para a falta de segurança nas implementações & Demonstração de ataques & EECSI 2017, pp.~1--6; \url{10.1109/EECSI.2017.8239179}; Só o resumo \\
\E{05} & Calabretta et al.~\cite{calabretta2018mqttauth} & C3 &
  MQTT-Auth: AugPAKE para confidencialidade e dois \emph{tokens} (autenticação e autorização por tópico); o \emph{broker} decifra e recifra (§IV-A) & Dois \emph{testbeds}: portáteis e UDOO Cortex-A9 (§V) & J.\ Commun.\ Softw.\ Syst.\ 14(4) (2018), pp.~320--331; \url{10.24138/jcomss.v14i4.604}; Editor (JCOMSS) \\
\E{06} & Hernández Ramos et al.~\cite{hernandez2018fuzzing} & C1 &
  Enquadramento de \emph{fuzzing} baseado em \emph{templates} para o MQTT & Resultados experimentais iniciais & Wirel.\ Commun.\ Mob.\ Comput.\ 2018, 8261746; \url{10.1155/2018/8261746}; Só o resumo \\
\E{07} & Dinculeană e Cheng~\cite{dinculeana2019vulnerabilities} & C5 &
  Mapeamento \emph{Value-to-HMAC}: confidencialidade e integridade com um só mecanismo; repetição fora do âmbito (§6) & Desempenho face a cifra simétrica em Onion Omega2+ (§3.1, §4.2) & Applied Sciences 9(5) (2019), 848; \url{10.3390/app9050848}; Editor (MDPI) \\
\E{08} & Malina et al.~\cite{malina2019secure} & C5 &
  Três níveis com assinatura Schnorr sobre nível, tópico, identidade, dados e carimbo temporal (§4.5, Eq.~6) & Analítica: contagem de operações (§5.1) & ARES 2019, pp.~1--10; \url{10.1145/3339252.3340503}; PDF em formato ACM (10~pp.) \\
\E{09} & Sadio et al.~\cite{sadio2019lightweight} & C5 &
  AEAD ChaCha20-Poly1305 fim-a-fim para MQTT e MQTT-SN & Arduino UNO; memória e tempo de execução (§IV-B) & IOTSMS 2019, pp.~119--123; \url{10.1109/IOTSMS48152.2019.8939177}; Editor (IEEE Xplore) \\
\E{10} & Vaccari et al.~\cite{vaccari2020slowite} & C1 &
  SlowITe: DoS de baixa taxa que explora a configuração do servidor pelo cliente & Serviços MQTT reais, com e sem cifra (§5) & Sensors 20(10) (2020), 2932; \url{10.3390/s20102932}; PubMed Central (PMC7285273) \\
\E{11} & Vaccari et al.~\cite{vaccari2020mqttset} & C7 &
  MQTTset: conjunto de dados de tráfego MQTT legítimo e de ataques & Sistema de deteção por ML (§3.2, §4) & Sensors 20(22) (2020), 6578; \url{10.3390/s20226578}; PubMed Central (PMC7698741) \\
\E{12} & Segarra et al.~\cite{segarra2020mqttz} & C6 &
  MQT-TZ: \emph{broker} em ARM TrustZone; TLS mútuo e cifra em duas camadas & Micro e macro \emph{benchmarks}; cargas de uma aplicação médica & SRDS 2020, pp.~256--265; \url{10.1109/SRDS51746.2020.00033}; Pré-publicação (arXiv:2007.12442v2) \\
\E{13} & Dahlmanns et al.~\cite{dahlmanns2021entrust} & C5 &
  ENTRUST: chave por tópico distribuída por um servidor de chaves com ACL; contador por publicador contra duplicados (§5.2--5.3) & RedBear Duo (Cortex-M3) e Mosquitto (§7) & ACM SaT-CPS 2021, pp.~78--87; \url{10.1145/3445969.3450423}; Cópia de autor \\
\E{14} & Spielvogel et al.~\cite{spielvogel2021tls} & C6 &
  O \emph{broker} só reencaminha por ligações com o nível de segurança exigido & Protótipo; medições de \emph{overhead} (§II-C) & STM 2021, LNCS, pp.~145--162; \url{10.1007/978-3-030-91859-0_8}; Pré-publicação (arXiv:2109.01169) \\
\E{15} & Prantl et al.~\cite{prantl2021performance} & C6 &
  Impacto do TLS~1.2 no MQTT: débito, estabelecimento de ligação, energia & ESP8266; todos os níveis de QoS (§4) & ACM/SPEC ICPE 2021, pp.~241--248; \url{10.1145/3427921.3450253}; Atas ICPE (SPEC Research) \\
\E{16} & Van Glabbeek et al.~\cite{vanglabbeek2022flexible} & C5 &
  Três opções com assinatura Schnorr e carimbo temporal; \emph{broker} honesto mas curioso, CA confiável (§3.1, §4.3) & $-80\,\%$ \emph{overhead} de comunicação e $-40\,\%$ computacional face ao MQTTS (Raspberry Pi) & Sensors 22(19) (2022), 7391; \url{10.3390/s22197391}; PubMed Central (PMC9572294) \\
\E{17} & Buccafurri et al.~\cite{buccafurri2023mqtta} & C8 &
  MQTT-A: anonimato de \emph{publishers} e subscritores via \emph{bridges} & \emph{Brokers} em nuvem em vários países; latência e \emph{goodput} (§VIII) & IEEE IoT J.\ 10(17) (2023), pp.~15443--15463; \url{10.1109/JIOT.2023.3264019}; Editor (IEEE Xplore) \\
\E{18} & Alotaibi et al.~\cite{alotaibi2024secure} & C7 &
  Deteção distribuída de anomalias com H2O (melhor: XGBoost) & MSE, RMSE, erro por classe, \emph{log loss} (§4.3) & Sensors 24(5) (2024), 1638; \url{10.3390/s24051638}; PubMed Central (PMC10935133) \\
\E{19} & Javadi et al.~\cite{javadi2026cmap} & C2 &
  CMAP: autenticação mútua sem certificados sobre ECC, entre dispositivo e \emph{broker} & Prova no modelo do oráculo aleatório (§V-B); Raspberry Pi com Mosquitto; NS-3 & IEEE Access 14 (2026), pp.~60836--60852; \url{10.1109/ACCESS.2026.3684371}; Editor (IEEE Access) \\
\midrule
--- & Perrone et al.~\cite{perrone2017mirai} & --- & Revisão relacionada (Secção~\ref{sec:relacionados}) & --- & IoTBDS 2017, pp.~246--253; \url{10.5220/0006287302460253}; Editor (SciTePress) \\
--- & Nebbione e Calzarossa~\cite{nebbione2020security} & --- & Revisão relacionada (Secção~\ref{sec:relacionados}) & --- & Future Internet 12(3) (2020), 55; \url{10.3390/fi12030055}; Editor (MDPI) \\
--- & Hintaw et al.~\cite{hintaw2023mqtt} & --- & Revisão relacionada (Secção~\ref{sec:relacionados}) & --- & IETE J.\ Res.\ 69(6) (2023), pp.~3368--3397; \url{10.1080/03772063.2021.1912651}; Só o resumo e a introdução \\
--- & Lakshminarayana et al.~\cite{lakshminarayana2024securing} & --- & Revisão relacionada (Secção~\ref{sec:relacionados}) & --- & IEEE Commun.\ Surv.\ Tutor.\ 26(4) (2024), pp.~2510--2546; \url{10.1109/COMST.2024.3372630}; Editor (IEEE Xplore) \\
--- & Laghari et al.~\cite{laghari2024securing} & --- & Revisão relacionada (Secção~\ref{sec:relacionados}) & --- & IEEE Access 12 (2024), pp.~139273--139289; \url{10.1109/ACCESS.2024.3412030}; Editor (IEEE Access) \\
\end{longtable}
}


\end{document}
